% Options for packages loaded elsewhere
\PassOptionsToPackage{unicode}{hyperref}
\PassOptionsToPackage{hyphens}{url}
\PassOptionsToPackage{dvipsnames,svgnames,x11names}{xcolor}
\documentclass[
  british,
  11pt,
  a4paper,
]{article}
\usepackage{xcolor}
\usepackage[margin=24mm]{geometry}
\usepackage{amsmath,amssymb}
\setcounter{secnumdepth}{-\maxdimen} % remove section numbering
\usepackage{iftex}
\ifPDFTeX
  \usepackage[T1]{fontenc}
  \usepackage[utf8]{inputenc}
  \usepackage{textcomp} % provide euro and other symbols
\else % if luatex or xetex
  \usepackage{unicode-math} % this also loads fontspec
  \defaultfontfeatures{Scale=MatchLowercase}
  \defaultfontfeatures[\rmfamily]{Ligatures=TeX,Scale=1}
\fi
\usepackage{lmodern}
\ifPDFTeX\else
  % xetex/luatex font selection
\fi
% Use upquote if available, for straight quotes in verbatim environments
\IfFileExists{upquote.sty}{\usepackage{upquote}}{}
\IfFileExists{microtype.sty}{% use microtype if available
  \usepackage[]{microtype}
  \UseMicrotypeSet[protrusion]{basicmath} % disable protrusion for tt fonts
}{}
\makeatletter
\@ifundefined{KOMAClassName}{% if non-KOMA class
  \IfFileExists{parskip.sty}{%
    \usepackage{parskip}
  }{% else
    \setlength{\parindent}{0pt}
    \setlength{\parskip}{6pt plus 2pt minus 1pt}}
}{% if KOMA class
  \KOMAoptions{parskip=half}}
\makeatother
\usepackage{longtable,booktabs,array}
\newcounter{none} % for unnumbered tables
\usepackage{calc} % for calculating minipage widths
% Correct order of tables after \paragraph or \subparagraph
\usepackage{etoolbox}
\makeatletter
\patchcmd\longtable{\par}{\if@noskipsec\mbox{}\fi\par}{}{}
\makeatother
% Allow footnotes in longtable head/foot
\IfFileExists{footnotehyper.sty}{\usepackage{footnotehyper}}{\usepackage{footnote}}
\makesavenoteenv{longtable}
\usepackage{graphicx}
\makeatletter
\newsavebox\pandoc@box
\newcommand*\pandocbounded[1]{% scales image to fit in text height/width
  \sbox\pandoc@box{#1}%
  \Gscale@div\@tempa{\textheight}{\dimexpr\ht\pandoc@box+\dp\pandoc@box\relax}%
  \Gscale@div\@tempb{\linewidth}{\wd\pandoc@box}%
  \ifdim\@tempb\p@<\@tempa\p@\let\@tempa\@tempb\fi% select the smaller of both
  \ifdim\@tempa\p@<\p@\scalebox{\@tempa}{\usebox\pandoc@box}%
  \else\usebox{\pandoc@box}%
  \fi%
}
% Set default figure placement to htbp
\def\fps@figure{htbp}
\makeatother
\ifLuaTeX
\usepackage[bidi=basic,shorthands=off]{babel}
\else
\usepackage[bidi=default,shorthands=off]{babel}
\fi
\ifLuaTeX
  \usepackage{selnolig} % disable illegal ligatures
\fi
\setlength{\emergencystretch}{3em} % prevent overfull lines
\providecommand{\tightlist}{%
  \setlength{\itemsep}{0pt}\setlength{\parskip}{0pt}}
\usepackage{amsmath,amssymb,mathtools,microtype,booktabs}
\providecommand{\mathscr}{\mathcal}
\DeclareMathOperator{\Var}{Var}
\DeclareMathOperator{\Cov}{Cov}
\DeclareMathOperator{\Rea}{Re}
\DeclareMathOperator{\Ran}{Ran}
\DeclareMathOperator{\diag}{diag}
\allowdisplaybreaks
\usepackage{bookmark}
\IfFileExists{xurl.sty}{\usepackage{xurl}}{} % add URL line breaks if available
\urlstyle{same}
\hypersetup{
  pdftitle={Sharp variance ordering for rotating Gaussian samplers},
  pdfauthor={Anonymous},
  pdflang={en-GB},
  colorlinks=true,
  linkcolor={black},
  filecolor={Maroon},
  citecolor={Blue},
  urlcolor={teal},
  pdfcreator={LaTeX via pandoc}}

\title{Sharp variance ordering for rotating Gaussian samplers}
\usepackage{etoolbox}
\makeatletter
\providecommand{\subtitle}[1]{% add subtitle to \maketitle
  \apptocmd{\@title}{\par {\large #1 \par}}{}{}
}
\makeatother
\subtitle{A matrix criterion, exact regularity and arithmetic
convergence rates}
\author{Anonymous}
\date{4 November 2026 --- research manuscript, version 1.4 --- DOI
\href{https://doi.org/10.5281/zenodo.23953250}{10.5281/zenodo.23953250}}

\begin{document}
\maketitle

\subsection{Abstract}\label{abstract}

Adding a target-preserving rotation to a kinetic Langevin sampler can
lower the long-run variance of time averages, but not always. Duncan,
Nüsken and Pavliotis (2017) proved a local improvement for linear and
quadratic observables of a Gaussian target when the friction exceeds
\(\sqrt2\), and conjectured the same second-derivative statement for
every square-integrable observable. We prove a global form of that
conjecture. With matched position and momentum rotations, every
square-integrable position observable has a long-run variance that never
increases with rotation strength if and only if the friction is at least
\(\sqrt2\). The literal second-derivative statement holds exactly on an
explicit weighted Hermite class, which contains every polynomial;
outside it we give observables whose variance has a
sharper-than-quadratic peak. For noncommuting stiffness \(K\) and
friction \(\Gamma\), the exact condition is
\(\Gamma^3-\Gamma K-K\Gamma\succeq0\), and the least friction at which
rotation is safe for every observable is \((2K)^{1/2}\), a safety
boundary rather than an optimum. Below the threshold, even observables
still improve at every rotation strength, and one finite rotation
strength improves every observable at once. We also give sharp
large-rotation rates and show which observation schedules preserve the
ordering. All results concern exact stationary Gaussian dynamics.

\textbf{Keywords:} kinetic Langevin dynamics; nonreversible sampling;
asymptotic variance; Gaussian chaos; symmetric unimodality; small
divisors.

\subsection{1. The question and the
contribution}\label{the-question-and-the-contribution}

\subsubsection{1.1. What is being
compared}\label{what-is-being-compared}

A Markov chain Monte Carlo sampler estimates an expectation by a time
average along one simulated trajectory. Its accuracy is measured by the
\emph{long-run} (asymptotic) variance of that average: the variance
after a long run of length \(T\), multiplied by \(T\). Two samplers with
the same target distribution can have very different long-run variances,
because successive states are correlated in different ways.

One way to alter those correlations without changing the target is to
add a deterministic rotation that preserves the target distribution. For
overdamped (first-order) Langevin samplers this is known to help every
observable: Hwang, Normand and Wu (2015) and Rey-Bellet and Spiliopoulos
(2015a, 2015b, 2016) show that such nonreversible perturbations never
increase the asymptotic variance, and that the variance decreases
monotonically with the perturbation strength. Duncan, Nüsken and
Pavliotis (2017), hereafter DNP, studied the analogous question for
\emph{kinetic} (second-order, underdamped) Langevin samplers, where the
position is driven by a momentum with friction \(\gamma\). They found
the kinetic case different: rotations can increase the variance.

\subsubsection{1.2. The question DNP left
open}\label{the-question-dnp-left-open}

DNP's Theorem 3 proves, for a Gaussian target with matched position and
momentum rotations of amplitude \(\mu\), that for linear and quadratic
observables, under nondegeneracy conditions, the variance has zero first
derivative and a negative second derivative at \(\mu=0\) when
\(\gamma>\sqrt2\). Immediately after the parameter conditions they
recommend (numbered (27) in the journal version and (35a)--(35c) in
arXiv:1705.00170v1), they conjecture that this derivative statement
holds for arbitrary square-integrable observables, and say they have not
been able to prove it. Nüsken's thesis (2018, p.~62) repeats the
conjecture. DNP's Theorem 5 identifies the large-amplitude limit as the
variance of the rotation-invariant part of the observable, and their
numerical experiments (Section 6) already show the variance rising for
small friction.

\subsubsection{1.3. What this paper
proves}\label{what-this-paper-proves}

\begin{quote}
\textbf{In plain terms.} Picture the position of the sampler as a damped
oscillator. Its \emph{power spectrum} says how much of its fluctuation
sits at each frequency. Rotating the sampler at speed \(\mu\) shifts
that spectrum, so the long-run variance becomes the spectrum evaluated
away from zero. If the spectrum is highest at zero frequency and falls
off on both sides, any shift lowers the variance. A damped oscillator
has such a single-peaked spectrum exactly when it shows no resonance
peak, and for the unit oscillator that is the familiar condition
\(\gamma\ge\sqrt2\) (quality factor at most \(1/\sqrt2\)). Nonlinear
observables see sums of convolution powers of the same spectrum, and
convolution keeps single-peaked shapes single-peaked. That is the whole
mechanism; the rest of the paper makes it exact, extends it to matrices,
and measures what happens at the edges.
\end{quote}

The main results are as follows.

{\def\LTcaptype{none} % do not increment counter
\begin{longtable}[]{@{}
  >{\raggedright\arraybackslash}p{(\linewidth - 4\tabcolsep) * \real{0.3333}}
  >{\raggedright\arraybackslash}p{(\linewidth - 4\tabcolsep) * \real{0.3333}}
  >{\raggedright\arraybackslash}p{(\linewidth - 4\tabcolsep) * \real{0.3333}}@{}}
\toprule\noalign{}
\begin{minipage}[b]{\linewidth}\raggedright
Result
\end{minipage} & \begin{minipage}[b]{\linewidth}\raggedright
Statement
\end{minipage} & \begin{minipage}[b]{\linewidth}\raggedright
Where
\end{minipage} \\
\midrule\noalign{}
\endhead
\bottomrule\noalign{}
\endlastfoot
Sharp ordering & Every \(L^2\) position observable has nonincreasing
variance in \(\lvert\mu\rvert\) \textbf{iff} \(\gamma\ge\sqrt2\); strict
unless rotation-invariant. The endpoint \(\gamma=\sqrt2\) is included. &
Theorem 1 \\
Spectral criterion & For any rotated stationary Gaussian process,
universal ordering holds iff the spectral density decreases on
\([0,\infty)\). & Theorem 2 \\
Below the threshold & For every \(\gamma>0\): even observables improve
at every \(\mu\ne0\); one finite amplitude improves all observables at
once. & Theorem 4A \\
Exact curvature domain & Every variance curve is \(C^1\); it is \(C^2\)
iff a weighted Hermite sum is finite. Explicit cusps outside. & Theorem
5, Prop. 6 \\
Large-rotation rates & \(O(\mu^{-4})\) for each polynomial; uniform
worst case \(\sim\sqrt{2\pi/e}/(\lvert\mu\rvert\omega_*)\); slower or no
algebraic rate for incommensurate speeds. & Prop. 7, Thm 8, §7 \\
Matrix criterion & For noncommuting \(K,\Gamma\): ordering iff
\(\Gamma^3-\Gamma K-K\Gamma\succeq0\); least friction \((2K)^{1/2}\);
two replicas make it apply to any Gaussian target. & Thm 10, Cors
11--12 \\
Observation schedules & Fixed horizons and regular sampling can reverse
the ordering; independent Poisson sampling preserves it. & Section 9 \\
\end{longtable}
}

Relative to DNP's conjecture, the precise outcome is this. The intended
content, that rotation reduces variance whenever the friction is large
enough, holds globally in \(\mu\) for every \(L^2\) position observable,
and also at the endpoint \(\gamma=\sqrt2\). The literal
second-derivative formulation needs two qualifications:
rotation-invariant observables have zero curvature (DNP's nondegeneracy
conditions exclude these), and some rough \(L^2\) observables have a
cusp, so that the second difference quotient tends to \(-\infty\). The
derivative statement holds exactly on an explicit class containing every
polynomial (Theorem 5).

\subsubsection{1.4. What is classical and what is
new}\label{what-is-classical-and-what-is-new}

Our method uses classical ingredients: Gaussian Hermite expansions,
Fourier transforms of covariance powers, and closure of symmetric
unimodality under convolution. Breuer and Major (1983, pp.~428, 430)
give the Hermite covariance and spectral-convolution identities; they
underlie the Breuer--Major central limit theorem. Maini (2026, Lemma 3.1
and Corollary 3.4) states the same convolution-power variance for
Fourier functionals of stationary Gaussian fields, one chaos at a time
and summed over chaoses, with a central limit theorem at each fixed
frequency; this is formula (6) below at a fixed frequency. That paper
does not compare spectral densities at different frequencies, and it has
no rotated family, Markov dynamics or sampler, so monotonicity in the
frequency does not arise there. Lefèvre and Utev (2013) discuss and
extend the convolution closure property. Gaussian spectral expansions
also appear in other applications, for example Muscinelli, Gerstner and
Schwalger (2019, equations (37), (42)--(43)). At critical damping, the
scalar spectral density belongs to the scaled Student family whose
convolution formulas are studied by Berg and Vignat (2008). We do not
claim these ingredients, or DNP's invariant-projection limit, as new.

The proposed contributions are the exact ordering criterion and kinetic
threshold, the below-threshold comparison, the regularity and rate
results, and the matrix criterion with its least friction. A related but
different line of work optimises nonreversible linear drifts and
friction for the \emph{rate of convergence} to equilibrium rather than
for variance ordering (Hwang, Hwang-Ma \& Sheu, 1993; Lelièvre, Nier \&
Pavliotis, 2013; Guillin \& Monmarché, 2016), or optimises friction for
a given observable (Chak et al., 2023); in discrete state spaces,
breaking detailed balance can also fail to help (Kaiser, Jack \& Zimmer,
2017). For drift-accelerated overdamped diffusions, Franke et al.~(2021)
prove that the variance tends to zero uniformly on the unit ball of
observables exactly when the drift operator has no nonzero \(H^1\)
kernel element (their Theorem 2.4 and Corollaries 3.1 and 4.3), and
their spectral decomposition of the variance is nonincreasing in the
drift amplitude for every observable. That argument needs a self-adjoint
unperturbed generator. The kinetic generator is not self-adjoint,
Theorem 1 shows that the ordering can then fail, and they give no rates,
so Theorems 1 and 8 below are different statements. Their introduction
notes that an irrational rotation of the 2-torus has no nontrivial
kernel; this is a qualitative example without Gaussian structure or
rates. Arithmetic effects on variance also appear in a different
Gaussian-field setting (Lachièze-Rey, 2022). Replica coupling and
replica averaging are established sampling devices (Section 8.4).

\subsubsection{1.5. How the paper is
organised}\label{how-the-paper-is-organised}

Section 2 states the model and the main theorem. Section 3 proves the
spectral criterion for general rotated Gaussian processes, and Section 4
applies it to the kinetic sampler, including the below-threshold
comparison (Section 4.1). Sections 5--7 concern regularity at small
rotations and rates at large rotations; readers interested only in the
ordering can skip them. Section 8 treats matrix stiffness and friction
and the replica construction. Section 9 explains which observation
schedules preserve the ordering. Section 10 describes the verification
scripts and the limits of the results. Appendix A collects explicit
checks and Appendix B the matrix versions of Sections 5--7.

\textbf{Notation.} \(V_f\) is the long-run variance (2); it is twice
DNP's \(\sigma_f^2\). \(S\), \(S_\gamma\) and \(S_n\) denote spectral
densities; the Gaussian target precision in Section 4 is written \(P\).
In Section 8, \(K\) and \(\Gamma\) are stiffness and friction matrices;
the Euler gamma function appears only in Proposition 6, written
\(\Gamma(\cdot)\) with a scalar argument.

\subsection{2. Model and main theorem}\label{model-and-main-theorem}

Let \(d\ge2\), let \(J\in\mathbb R^{d\times d}\) satisfy \(J^\top=-J\),
and let \(\gamma>0\). We consider

\[
\begin{aligned}
\mathrm dq_t&=(p_t-\mu Jq_t)\,\mathrm dt,\\
\mathrm dp_t&=(-q_t-\gamma p_t-\mu Jp_t)\,\mathrm dt
                 +\sqrt{2\gamma}\,\mathrm dW_t .
\end{aligned} \tag{1}
\]

Time and coordinates are normalised so that the harmonic frequency and
the position and momentum variances are one. The stationary law is
\(\widehat\pi=\mathcal N(0,I_{2d})\), and the position marginal is
\(\pi=\mathcal N(0,I_d)\). We start the process in stationarity. An
observable is a real function \(f\in L^2(\pi)\) of position alone.
Subtracting its mean changes none of the results, so we take
\(\pi(f)=0\).

Define its long-run variance by

\[
V_f(\mu)=\lim_{T\to\infty}T\mathop{\mathrm{Var}}\left(
\frac1T\int_0^T f(q_t)\,\mathrm dt\right). \tag{2}
\]

We prove existence of this limit directly. This convention is twice
DNP's \(\sigma_f^2\). We use the variance limit as the definition; no
central limit theorem is needed in our proofs.

Write \(U_\theta f(q)=f(e^{-\theta J}q)\). Let \(\Pi_J\) denote the
orthogonal projection in \(L^2(\pi)\) onto the functions invariant under
every \(U_\theta\). It preserves centred functions and fixes constants.

\textbf{Theorem 1 (sharp global ordering).} Suppose \(J\ne0\). The
following are equivalent:

\begin{enumerate}
\def\labelenumi{\arabic{enumi}.}
\tightlist
\item
  \(\gamma\ge\sqrt2\).
\item
  For every real \(f\in L^2(\pi)\) and every \(0\le a<b\),
  \(V_f(b)\le V_f(a)\).
\end{enumerate}

When these conditions hold, equality at any one such pair occurs if and
only if \(f=\Pi_Jf\). Thus each non-invariant observable has strictly
decreasing variance on \([0,\infty)\). The variance is even in \(\mu\).
If \(0<\gamma<\sqrt2\), a linear observable has \(V_f(\mu)>V_f(0)\) for
all sufficiently small nonzero \(\mu\).

The threshold includes \(\gamma=\sqrt2\), although a linear observable
then has zero second curvature at the origin. If \(J=0\), all variances
are constant for every \(\gamma\), so the nonzero assumption is
necessary for the equivalence.

Theorem 1 concerns the Gaussian model with exactly matched rotations. It
does not assert the same ordering for general target distributions,
arbitrary phase-space observables, or unmatched perturbations. Section 9
gives explicit observation-related boundaries.

\subsection{3. A spectral characterisation for Gaussian
processes}\label{a-spectral-characterisation-for-gaussian-processes}

We first prove a statement that does not require a Markov process.

\textbf{Theorem 2 (spectral criterion).} Let \((X_t)_{t\in\mathbb R}\)
be a centred, mean-square continuous, stationary Gaussian process in
\(\mathbb R^d\) satisfying

\[
\mathop{\mathrm{Cov}}(X_t,X_s)=r(t-s)I_d,\qquad r(0)=1,\qquad r\in L^1(\mathbb R).
\]

Fix a nonzero real skew matrix \(J\), and put
\(Q_t^\mu=e^{-\mu tJ}X_t\). Use the Fourier convention

\[
S(\xi)=\widehat r(\xi)=\int_{\mathbb R}e^{-i\xi t}r(t)\,\mathrm dt .
\]

Every real centred \(f\in L^2(\pi)\) has a finite variance limit (2),
with \(q\) replaced by \(Q\). These limits are nonincreasing in
\(|\mu|\) for every such \(f\) if and only if \(S\) is nonincreasing on
\([0,\infty)\). If \(S\) is strictly decreasing there, the ordering is
strict precisely for non-invariant observables.

\subsubsection{3.1. Hermite and rotation
decomposition}\label{hermite-and-rotation-decomposition}

The process \(Q^\mu\) is stationary because it is Gaussian with constant
mean and covariance

\[
\mathop{\mathrm{Cov}}(Q_t^\mu,Q_s^\mu)=r(t-s)e^{-\mu(t-s)J}.
\]

Write \(L^2_0(\pi)=\bigoplus_{n\ge1}\mathcal H_n\) for the orthogonal
decomposition into position Hermite chaoses. Here \(\mathcal H_n\) is
the span of multivariate Hermite polynomials of total degree \(n\). Each
space is finite dimensional. The Mehler operator

\[
T_\rho f(x)=\mathbb E\big[f(\rho x+\sqrt{1-\rho^2}Z)\big],
\qquad Z\sim\pi,
\]

acts as multiplication by \(\rho^n\) on \(\mathcal H_n\), also for
negative \(\rho\). Orthogonal rotations preserve each chaos and commute
with \(T_\rho\).

Complexify each chaos and diagonalise the unitary group \(U_\theta\).
Let \(E_{n,\omega}\) be its spectral projection with

\[
U_\theta E_{n,\omega}=e^{i\omega\theta}E_{n,\omega},
\qquad
a_{n,\omega}=\|E_{n,\omega}f_n\|_2^2.
\]

The index \(\omega\) ranges over the finite set of rotation frequencies
in degree \(n\). These weights satisfy

\[
a_{n,\omega}\ge0,\quad a_{n,-\omega}=a_{n,\omega},\quad
\sum_{n,\omega}a_{n,\omega}=\|f\|_2^2,
\quad |\omega|\le n\|J\|. \tag{3}
\]

Both signs of a nonzero frequency are included in this convention; their
weights are not counted twice elsewhere.

Joint Gaussian conditioning gives

\[
\mathbb E[f(Q_t^\mu)\mid Q_0^\mu=x]
=T_{r(t)}U_{\mu t}f(x).
\]

Orthogonality and (3) therefore give the covariance expansion

\[
C_f(t;\mu)=\sum_{n,\omega}a_{n,\omega}r(t)^n\cos(\mu\omega t). \tag{4}
\]

As a covariance normalised at zero, \(|r(t)|\le1\). Consequently,

\[
|C_f(t;\mu)|\le\|f\|_2^2|r(t)|,
\qquad
\sum_{n,\omega}a_{n,\omega}\int_{\mathbb R}|r(t)|^n\,\mathrm dt
\le\|f\|_2^2\|r\|_1. \tag{5}
\]

These bounds justify all subsequent exchanges of sums and integrals. The
stationary variance identity

\[
T\mathop{\mathrm{Var}}(\overline f_T)=2\int_0^T(1-t/T)C_f(t;\mu)\,\mathrm dt
\]

and dominated convergence prove existence of (2).

Since \(r\) is continuous, positive definite and integrable, \(S\) is a
continuous nonnegative integrable spectral density and \(\int S=2\pi\).
Set \(p=S/(2\pi)\) and

\[
S_n(\xi)=\widehat{r^n}(\xi)
=(2\pi)^{1-n}S^{*n}(\xi)=2\pi p^{*n}(\xi).
\]

We obtain the exact nonnegative expansion

\[
\boxed{\displaystyle V_f(\mu)=\sum_{n\ge1,\omega}
a_{n,\omega}S_n(\mu\omega).} \tag{6}
\]

Time-domain covariances can change sign. Their spectral densities in (6)
remain nonnegative.

\subsubsection{3.2. Convolution lemma and proof of the
criterion}\label{convolution-lemma-and-proof-of-the-criterion}

\textbf{Lemma 3.} The convolution of two even, nonnegative, integrable
functions that are nonincreasing on \([0,\infty)\) has the same
properties. If one is continuous and strictly decreasing on
\([0,\infty)\), and the other has strictly positive integral, their
convolution is strictly decreasing on \([0,\infty)\).

\textbf{Proof.} Represent the second function by its superlevel
intervals, using the layer-cake identity. For each interval radius
\(R>0\), convolution of the first function \(a\) with its indicator
gives

\[
H_R(x)=\int_{x-R}^{x+R}a(u)\,\mathrm du .
\]

If \(a\) is continuous and strictly decreasing away from zero, then

\[
H_R'(x)=a(x+R)-a(x-R)<0\qquad(x>0),
\]

because \(a\) is even and \(|x-R|<x+R\). The weak version needs no
differentiability: for \(0\le y<x\),

\[
H_R(x)-H_R(y)=\int_{y+R}^{x+R}a(u)\,\mathrm du-\int_{y-R}^{x-R}a(u)\,\mathrm du\le0,
\]

because the two intervals have equal length and, as \(a\) is even and
nonincreasing in \(|u|\), each point \(u+2R\) of the first is at least
as far from the origin as the corresponding point \(u\) of the second.
Integrate over the superlevel intervals. When the second function has
strictly positive integral, positive radii occur on a set of levels of
positive measure, which preserves strictness between any two distinct
nonnegative arguments. \(\square\)

If \(S\) is nonincreasing, Lemma 3 applies to every convolution power in
(6), proving sufficiency in Theorem 2. Under strict decrease, every
summand with \(\omega\ne0\) and positive weight decreases strictly. Such
a summand exists exactly when \(f\ne\Pi_Jf\).

For necessity, choose a two-dimensional rotation plane of \(J\) with
nonzero speed \(\lambda>0\), and a unit linear observable in that plane.
Its covariance is \(r(t)\cos(\lambda\mu t)\) and its variance is
\(S(\lambda\mu)\). As \(\mu\) traverses \([0,\infty)\), so does
\(\lambda\mu\). Universal monotonicity therefore forces monotonicity of
\(S\). \(\square\)

\subsection{4. Applying the criterion to kinetic Langevin
dynamics}\label{applying-the-criterion-to-kinetic-langevin-dynamics}

The drift matrix of (1) is

\[
B_\mu=\begin{pmatrix}\mu J&-I\\I&\gamma I+\mu J\end{pmatrix}
=B_0+\mu\mathop{\mathrm{diag}}(J,J).
\]

The two summands commute. Also
\(B_\mu+B_\mu^\top=2\mathop{\mathrm{diag}}(0,\gamma I)\), and all drift
eigenvalues have positive real parts. Thus \(I_{2d}\) is the stationary
covariance, and

\[
\mathop{\mathrm{Cov}}(q_t,q_0)=r_\gamma(t)e^{-\mu tJ},
\quad r_\gamma''+\gamma r_\gamma'+r_\gamma=0,
\quad r_\gamma(0)=1,\quad r_\gamma'(0)=0. \tag{7}
\]

The scalar covariance decays exponentially, allowing a polynomial
prefactor at \(\gamma=2\). Its Laplace transform is

\[
R_\gamma(s)=\frac{s+\gamma}{s^2+\gamma s+1}.
\]

Extending \(r_\gamma\) evenly, we find

\[
S_\gamma(\xi)=2\mathop{\mathrm{Re}}R_\gamma(i\xi)
=\frac{2\gamma}{\xi^4+(\gamma^2-2)\xi^2+1}, \tag{8}
\]

and

\[
S_\gamma'(\xi)=
\frac{-4\gamma\xi(2\xi^2+\gamma^2-2)}
{[\xi^4+(\gamma^2-2)\xi^2+1]^2}. \tag{9}
\]

For \(\gamma\ge\sqrt2\), (9) is strictly negative for every \(\xi>0\).
Theorem 2 proves the positive direction of Theorem 1, including its
equality statement.

Conversely, the unit linear observable used in the necessity argument
has

\[
V_f(\mu)=\frac{2\gamma}
{1+(\gamma^2-2)(\lambda\mu)^2+(\lambda\mu)^4}. \tag{10}
\]

For \(\gamma<\sqrt2\), this exceeds \(V_f(0)=2\gamma\) whenever
\(0<(\lambda\mu)^2<2-\gamma^2\). This proves the converse, and is the
mechanism behind the increase that DNP observed numerically at small
friction. At \(\gamma=\sqrt2\), its leading decrease is fourth order,
which explains why global strictness and negative second curvature have
different endpoint conditions. \(\square\)

\textbf{General Gaussian targets.} Let the target have precision matrix
\(P>0\). In DNP's notation, take mass \(M=P\), friction
\(\Gamma=\gamma P\), \(J_2=PJ_1P\), and equal perturbation amplitudes;
this is exactly DNP's recommended choice (their conditions (27) in the
journal version). The change of variables

\[
\widetilde q=P^{1/2}q,\quad
\widetilde p=P^{-1/2}p,\quad
\widetilde J=P^{1/2}J_1P^{1/2}
\]

reduces the dynamics to (1). Apply Theorem 1 to
\(f(P^{-1/2}\widetilde q)\). This statement requires the specified
preconditioning; it gives no uniform tolerance to covariance
misspecification.

\subsubsection{4.1. Below the threshold: comparison with the unrotated
sampler}\label{below-the-threshold-comparison-with-the-unrotated-sampler}

Theorem 1 says that for \(\gamma<\sqrt2\) some observable gets worse at
some small rotation. Practitioners often use small friction, so the
useful question is weaker: which observables, and which amplitudes,
still improve on the \emph{unrotated} sampler? Here we compare
\(V_f(\mu)\) with \(V_f(0)\) rather than asking for monotonicity in
\(|\mu|\). For each \(n\ge1\) define the crossing point

\[
\Xi_n(\gamma)=\sup\{\xi>0:\ S_n(\xi)\ge S_n(0)\},\qquad\sup\varnothing=0 .
\]

By (8), \(\Xi_1(\gamma)=\sqrt{2-\gamma^2}\) for \(\gamma<\sqrt2\) and
\(\Xi_1=0\) otherwise.

\textbf{Theorem 4A (below-threshold comparison).} Let \(\gamma>0\) and
\(J\ne0\).

\begin{enumerate}
\def\labelenumi{\arabic{enumi}.}
\tightlist
\item
  \emph{Even observables.} If \(f\) has no component in odd Hermite
  chaoses (for instance, if \(f(-q)=f(q)\)), then \(V_f(\mu)<V_f(0)\)
  for every \(\mu\ne0\) unless \(f=\Pi_Jf\).
\item
  \emph{Only finitely many degrees can be hurt.} There is an integer
  \(N(\gamma)\ge1\) such that \(S_n(\xi)<S_n(0)\) for every \(\xi\ne0\)
  and every \(n>N(\gamma)\). Every \(\Xi_n(\gamma)\) is finite.
\item
  \emph{One amplitude for all observables.} Let \(\omega_{\min}>0\) be
  the smallest nonzero rotation frequency occurring in Hermite degrees
  at most \(N(\gamma)\), and set \[
  \mu_0(\gamma,J)=\max_{1\le n\le N(\gamma)}\frac{\Xi_n(\gamma)}{\omega_{\min}} .
  \] Then \(V_f(\mu)<V_f(0)\) for every real \(f\in L^2(\pi)\) with
  \(f\ne\Pi_Jf\) and every \(|\mu|>\mu_0(\gamma,J)\).
\item
  \emph{The amplitude cannot be small.} If \(\gamma<\sqrt2\) and
  \(\lambda_{\min}\) is the smallest positive rotation speed of \(J\), a
  linear observable in that rotation plane has \(V_f(\mu)>V_f(0)\) for
  \(0<|\mu|<\sqrt{2-\gamma^2}/\lambda_{\min}\).
\end{enumerate}

Part 1 does not assert monotonicity: for \(\gamma<\sqrt2\) even chaoses
can have non-monotone spectra. When the rotation speeds are commensurate
with gap \(\omega_*\), part 3 holds with \(\omega_{\min}\) replaced by
\(\omega_*\).

\textbf{Proof.} By (6),
\(V_f(0)-V_f(\mu)=\sum_{n,\omega}a_{n,\omega}[S_n(0)-S_n(\mu\omega)]\),
and

\[
S_n(0)-S_n(\xi)=2\int_0^\infty r_\gamma(t)^n\,(1-\cos\xi t)\,\mathrm dt. \tag{10a}
\]

\emph{Part 1.} For even \(n\), \(r_\gamma^n\ge0\) and is positive near
\(t=0\), so (10a) is strictly positive for every \(\xi\ne0\). Every term
with \(\omega\ne0\) and positive weight therefore contributes strictly,
and such a term exists exactly when \(f\ne\Pi_Jf\).

\emph{Part 2.} Take \(\delta>0\), \(C>0\) and \(\rho<1\) with
\(r_\gamma(t)\ge e^{-Ct^2}\) on \([0,\delta]\) and
\(|r_\gamma(t)|\le\rho\) for \(t\ge\delta\); these exist by the
expansion and energy argument in the proof of Lemma 4. On
\([\delta,\infty)\) bound the integrand of (10a) below by
\(-|r_\gamma|^n(1-\cos\xi t)\) and use
\(1-\cos\xi t\le(2+t^2/2)\min(1,\xi^2)\). This gives

\[
S_n(0)-S_n(\xi)\ge 2\int_0^\delta e^{-Cnt^2}(1-\cos\xi t)\,\mathrm dt-2M_\gamma\rho^{n-1}\min(1,\xi^2),
\qquad M_\gamma=\int_\delta^\infty|r_\gamma(t)|\bigl(2+\tfrac{t^2}2\bigr)\mathrm dt .
\]

We bound the first integral below by \(c\,n^{-3/2}\min(1,\xi^2)\) for
all large \(n\). If \(|\xi|\le\pi/\delta\), then
\(1-\cos u\ge2u^2/\pi^2\) on \(|u|\le\pi\) gives at least
\((2\xi^2/\pi^2)\,n^{-3/2}\int_0^{\delta\sqrt n}u^2e^{-Cu^2}\,\mathrm du\),
and the last integral is bounded below for \(n\ge1\). If
\(|\xi|>\pi/\delta\), write the integral as
\(\int_0^\delta e^{-Cnt^2}\mathrm dt-\int_0^\delta e^{-Cnt^2}\cos\xi t\,\mathrm dt\).
Up to errors \(O(e^{-Cn\delta^2})\), these are
\(\tfrac12\sqrt{\pi/(Cn)}\) and
\(\tfrac12\sqrt{\pi/(Cn)}\,e^{-\xi^2/(4Cn)}\), so the difference is at
least
\(\tfrac12\sqrt{\pi/(Cn)}\min\{\tfrac12,\pi^2/(8C\delta^2n)\}-O(e^{-Cn\delta^2})\),
which is at least \(c\,n^{-3/2}\) for large \(n\). Since \(n^{-3/2}\)
dominates \(\rho^{n-1}\), there is \(N(\gamma)\) beyond which (10a) is
strictly positive for every \(\xi\ne0\).

For \(n\le N(\gamma)\), \(S_n=2\pi p^{*n}\) is a convolution power of
the everywhere-positive density \(p=S_\gamma/(2\pi)\), so \(S_n(0)>0\),
while \(S_n(\xi)\to0\) as \(|\xi|\to\infty\) by the Riemann--Lebesgue
lemma. Hence \(\Xi_n(\gamma)<\infty\).

\emph{Part 3.} Each degree contains finitely many frequencies, so
\(\omega_{\min}>0\). Let \(|\mu|>\mu_0\) and \(\omega\ne0\) be a
frequency in degree \(n\). If \(n>N(\gamma)\), part 2 gives
\(S_n(\mu\omega)<S_n(0)\). If \(n\le N(\gamma)\), then
\(|\mu\omega|\ge|\mu|\omega_{\min}>\Xi_n(\gamma)\), and
\(S_n(\mu\omega)<S_n(0)\) by definition of \(\Xi_n\). Summing over (6)
proves part 3.

\emph{Part 4} is (10) with \(\lambda=\lambda_{\min}\). \(\square\)

\textbf{How large is \(N(\gamma)\)?} The proof gives existence, not a
value. Table 1 reports the crossing points computed by
\texttt{code/verify\_below\_threshold.py}, from an exact exponential-sum
form of \(S_n\) (cross-checked by quadrature), scanning odd degrees up
to \(201\). These are numerical computations, not part of the proof.
Even degrees never exceed \(S_n(0)\), as part 1 requires.

Table 1. Crossing points \(\Xi_n(\gamma)\) for odd degrees \(n\)
(computed, odd \(n\le121\)).

{\def\LTcaptype{none} % do not increment counter
\begin{longtable}[]{@{}
  >{\raggedright\arraybackslash}p{(\linewidth - 6\tabcolsep) * \real{0.2500}}
  >{\raggedright\arraybackslash}p{(\linewidth - 6\tabcolsep) * \real{0.2500}}
  >{\raggedright\arraybackslash}p{(\linewidth - 6\tabcolsep) * \real{0.2500}}
  >{\raggedright\arraybackslash}p{(\linewidth - 6\tabcolsep) * \real{0.2500}}@{}}
\toprule\noalign{}
\begin{minipage}[b]{\linewidth}\raggedright
\(\gamma\)
\end{minipage} & \begin{minipage}[b]{\linewidth}\raggedright
last odd \(n\) with a crossing
\end{minipage} & \begin{minipage}[b]{\linewidth}\raggedright
\(\max_n\Xi_n\)
\end{minipage} & \begin{minipage}[b]{\linewidth}\raggedright
\(n\): \(\Xi_n\)
\end{minipage} \\
\midrule\noalign{}
\endhead
\bottomrule\noalign{}
\endlastfoot
0.1 & 41 & 3.358 & 1: 1.411, 3: 3.258, 5: 3.331, 7: 3.358, \ldots, 41:
0.551 \\
0.3 & 11 & 1.489 & 1: 1.382, 3: 1.489, 5: 1.427, 7: 1.270, 9: 0.996, 11:
0.426 \\
0.5 & 5 & 1.323 & 1: 1.323, 3: 1.257, 5: 0.784 \\
0.8 & 3 & 1.166 & 1: 1.166, 3: 0.047 \\
1.0 & 1 & 1.000 & 1: 1.000 \\
1.2 & 1 & 0.748 & 1: 0.748 \\
\end{longtable}
}

Two features matter in practice. First, for moderate friction (in this
scan, \(\gamma=0.8\) and above) only the first one or two odd chaoses
can be hurt, and for \(\gamma\ge1\) only the linear component: then
\(\mu_0=\sqrt{2-\gamma^2}/\lambda_{\min}\) is attained by part 4.
Second, at small friction the set of harmful amplitudes is not an
interval, and \(N(\gamma)\) grows roughly like \(1/\gamma\). The
spectrum \(S_\gamma\) then has a sharp resonance peak near the damped
frequency \(\sqrt{1-\gamma^2/4}\), and odd chaoses carry peaks near its
odd multiples. A rotation that moves weight onto such a peak raises the
variance.

\textbf{Friction versus rotation.} The ordering threshold is a safety
boundary, not an optimum. For a unit linear observable, (10) gives
\(V_f(0)=2\gamma\) without rotation, which tends to zero as
\(\gamma\to0\); DNP found their best results at small friction. Theorem
1 and Corollary 11 below say when rotation can be added with no risk to
any observable, not which friction minimises a particular variance. The
threshold \(\gamma=\sqrt2\) also lies below critical damping
\(\gamma=2\), which maximises the spectral gap of the unperturbed scalar
oscillator.

\subsection{5. The exact regularity
boundary}\label{the-exact-regularity-boundary}

In this section \(\gamma>0\) is arbitrary and
\(S_n=\widehat{r_\gamma^n}\). Thus the regularity statements also apply
below the monotonicity threshold.

\subsubsection{5.1. Moment and Fourier
estimates}\label{moment-and-fourier-estimates}

\textbf{Lemma 4.} For every nonnegative integer \(k\), there is a finite
\(C_{\gamma,k}\) such that

\[
\int_0^\infty t^k|r_\gamma(t)|^n\,\mathrm dt
\le C_{\gamma,k}n^{-(k+1)/2}. \tag{11}
\]

Moreover,

\[
n^{(k+1)/2}\int_0^\infty t^kr_\gamma(t)^n\,\mathrm dt
\longrightarrow\int_0^\infty u^ke^{-u^2/2}\,\mathrm du. \tag{12}
\]

In particular,

\[
-S_n''(0)\sim\sqrt{2\pi}\,n^{-3/2},\qquad
\sqrt n S_n(\sqrt n x)\longrightarrow\sqrt{2\pi}e^{-x^2/2}, \tag{13}
\]

the latter uniformly over \(x\in\mathbb R\).

\textbf{Proof.} The oscillator energy has derivative
\((r^2+(r')^2)'=-2\gamma(r')^2\). Since \(r''(0)=-1\), its energy is
strictly below one at every positive time. Hence \(|r(t)|<1\) for
\(t>0\). At zero,

\[
r(t)=1-\tfrac12t^2+\tfrac\gamma6t^3+O(t^4).
\]

Choose \(\delta>0\) so that \(r(t)>0\) and \(r(t)\le e^{-ct^2}\) on
\([0,\delta]\). On the complement, \(\rho=\sup_{t\ge\delta}|r(t)|<1\),
and

\[
\int_\delta^\infty t^k|r(t)|^n\,\mathrm dt
\le\rho^{n-1}\int_\delta^\infty t^k|r(t)|\,\mathrm dt .
\]

The small-time Gaussian bound proves (11). Substitution \(u=\sqrt n\,t\)
and dominated convergence prove (12). The even rescaled functions
\(r(|u|/\sqrt n)^n\) converge in \(L^1(\mathrm du)\) to \(e^{-u^2/2}\)
by the same split; Fourier transformation proves the uniform assertion
in (13). \(\square\)

\subsubsection{5.2. First and second
derivatives}\label{first-and-second-derivatives}

\textbf{Theorem 5 (exact curvature domain).} Every real
\(f\in L^2(\pi)\) has \(V_f\in C^1(\mathbb R)\) and \(V_f'(0)=0\).
Define

\[
\mathcal Q_J(f)=\sum_{n,\omega}a_{n,\omega}\omega^2 n^{-3/2}. \tag{14}
\]

The following are equivalent: \(\mathcal Q_J(f)<\infty\); \(V_f\) has a
finite ordinary second derivative at zero; and
\(V_f\in C^2(\mathbb R)\). Under these conditions,

\[
V_f''(\mu)=-2\sum_{n,\omega}a_{n,\omega}\omega^2
\int_0^\infty t^2r_\gamma(t)^n\cos(\mu\omega t)\,\mathrm dt . \tag{15}
\]

If \(\mathcal Q_J(f)=\infty\), then

\[
\frac{V_f(\mu)-V_f(0)}{\mu^2}\longrightarrow-\infty
\quad\text{as }\mu\to0. \tag{16}
\]

Equation (16) concerns the symmetric difference quotient. It does not,
by itself, assert an extended derivative of \(V_f'\).

\textbf{Proof.} By (3) and (11),

\[
|\omega S_n'(\mu\omega)|
\le2|\omega|\int_0^\infty t|r(t)|^n\,\mathrm dt
\le2\|J\|C_{\gamma,1}.
\]

The first-derivative series in (6) therefore converges uniformly,
because the weights sum to \(\|f\|_2^2\). Each term has derivative zero
at zero. This proves the \(C^1\) assertion. If (14) is finite, the bound
\(|\omega^2\int_0^\infty t^2r(t)^n\cos(\mu\omega t)\,\mathrm dt|\le C_{\gamma,2}\,\omega^2n^{-3/2}\),
which is (11) with \(k=2\), gives uniform convergence of the
second-derivative series on \(\mathbb R\). This proves sufficiency,
continuity of \(V_f''\) and (15).

For necessity, the possibility that \(r(t)\) is negative must be
handled. Decompose \(r(t)^n=P_n(t)-N_n(t)\) into its nonnegative
positive and negative parts. Every negative part lies in
\([\delta,\infty)\) for the \(\delta\) in Lemma 4, and

\[
\int_0^\infty t^2N_n(t)\,\mathrm dt\le B_\gamma\rho^{n-1}.
\]

Using \(1-\cos x\le x^2/2\), its total contribution to

\[
\frac{V_f(0)-V_f(\mu)}{\mu^2}
=2\sum_{n,\omega}a_{n,\omega}
\int_0^\infty r(t)^n\frac{1-\cos(\mu\omega t)}{\mu^2}\,\mathrm dt
\]

is bounded uniformly in \(\mu\ne0\) by

\[
B_\gamma\|J\|^2
\sup_{n\ge1}(n^2\rho^{n-1})\,\|f\|_2^2<\infty.
\]

The positive-part second moments are asymptotic to
\(\sqrt{\pi/2}\,n^{-3/2}\) by (12). If (14) diverges, Fatou's lemma
applied to the positive-part integrands makes their contribution diverge
to \(+\infty\). The bounded negative contribution cannot cancel it. This
proves (16). A finite ordinary second derivative, together with
\(V_f'(0)=0\), would give a finite symmetric difference quotient, so
(14) is necessary. \(\square\)

In spectral notation, \(\mathcal Q_J(f)=\|N^{-3/4}R_Jf\|_2^2\), where
\(N\) is the positive Hermite number operator and
\(R_J=(Jq)\cdot\nabla\) is interpreted spectrally. A sufficient
condition is \(f\in D(N^{1/4})\), since \(|\omega|\le n\|J\|\). It is
not necessary: invariant components require no extra regularity.

\textbf{Curvature sign above the threshold.} For \(\gamma>\sqrt2\),
every \(S_n''(0)<0\). For \(n=1\) this follows from (8). For \(n\ge2\),
write \(p=S_\gamma/(2\pi)\) and \(h=p^{*(n-1)}\). Their smooth
derivatives are negative on \((0,\infty)\). To verify this statement
inductively, if \(g\) is an even smooth density with \(g'(y)<0\) for
\(y>0\), then

\[
(p*g)'(x)=\int_0^\infty g'(y)\,[p(x-y)-p(x+y)]\,\mathrm dy<0
\qquad(x>0).
\]

The bracket is positive because \(|x-y|<x+y\) and \(p\) is even and
strictly decreasing. Integration by parts now gives

\[
(p*h)''(0)=-\int_{\mathbb R}p'(y)h'(y)\,\mathrm dy<0.
\]

The integrability and boundary conditions follow from the rational decay
of \(p\) and its derivatives and their convolution representations. The
same argument applies for \(n\ge2\) when \(\gamma=\sqrt2\); only \(n=1\)
then has zero curvature. Thus, for \(\gamma>\sqrt2\), the derivative
formulation conjectured by DNP is valid precisely on the
finite-curvature domain (14), with strict negativity for non-invariant
\(f\). Global strict ordering needs only \(f\in L^2\).

\subsubsection{5.3. Explicit fractional and logarithmic
cusps}\label{explicit-fractional-and-logarithmic-cusps}

Take \(d=2\) and \(J=\begin{psmallmatrix}0&-1\\1&0\end{psmallmatrix}\).
Set

\[
\phi_n(q)=\frac{\mathop{\mathrm{Re}}(q_1+iq_2)^n}{\sqrt{2^{n-1}n!}},\qquad n\ge1.
\]

These are real orthonormal chaos elements. The polynomial is homogeneous
and harmonic, so its Hermite degree is \(n\); rotational symmetry and
\(\mathbb E|q_1+iq_2|^{2n}=2^n n!\) give the normalisation. Its angular
frequencies are \(\pm n\), and its covariance is \(r(t)^n\cos(n\mu t)\).

For \(1<p\le3/2\), define the \(L^2\)-convergent series

\[
f_p=\sum_{m\ge1}m^{-p/2}\phi_{2m},\qquad
\|f_p\|_2^2=\zeta(p). \tag{17}
\]

\textbf{Proposition 6.} For every \(\gamma>0\) and \(1<p<3/2\),

\[
V_{f_p}(0)-V_{f_p}(\mu)
\sim\frac{\sqrt\pi\,\Gamma(3/2-p)}{p-1/2}|\mu|^{2p-1}, \tag{18}
\]

and

\[
V_{f_p}'(\mu)\sim-2\sqrt\pi\,\Gamma(3/2-p)
\operatorname{sgn}(\mu)|\mu|^{2p-2}. \tag{19}
\]

At \(p=3/2\),

\[
\begin{aligned}
V_{f_{3/2}}(0)-V_{f_{3/2}}(\mu)
&\sim2\sqrt\pi\,\mu^2\log(1/|\mu|),\\
V_{f_{3/2}}'(\mu)&\sim-4\sqrt\pi\,\mu\log(1/|\mu|).
\end{aligned} \tag{20}
\]

\textbf{Proof.} Formula (6) gives
\(V_{f_p}(\mu)=\sum m^{-p}S_{2m}(2m\mu)\). The even powers of \(r\) make
all time kernels nonnegative. With \(h=\mu^2\) and \(x=mh\), (13) gives
the locally uniform scaled difference

\[
\frac{S_{2m}(0)-S_{2m}(2m\mu)}{|\mu|}
\longrightarrow\sqrt{\pi/x}(1-e^{-x}).
\]

Lemma 4 and the two elementary cosine bounds give

\[
0\le S_{2m}(0)-S_{2m}(2m\mu)
\le C\min\{m^{-1/2},\mu^2m^{1/2}\}. \tag{21}
\]

After rescaling, the summands are dominated by a constant times
\(\min(x^{1/2-p},x^{-p-1/2})\). Its integral is finite at zero and
infinity for the stated range. Splitting the Riemann sums into a compact
interval and two tails yields

\[
\lim_{\mu\to0}\frac{V_{f_p}(0)-V_{f_p}(\mu)}{|\mu|^{2p-1}}
=\sqrt\pi\int_0^\infty x^{-p-1/2}(1-e^{-x})\,\mathrm dx.
\]

Integration by parts evaluates this integral as the constant in (18).
The derivative scaling follows by differentiating the Fourier integral
after the \(u=\sqrt n t\) substitution. In particular,
\(nS_n'(\sqrt n x)\to-\sqrt{2\pi}x e^{-x^2/2}\) on compact sets. The
bound

\[
|2mS_{2m}'(2m\mu)|\le C\min\{1,|\mu|\sqrt m\}
\]

supplies the integrable Riemann-sum envelope \(\min(x^{-p},x^{1/2-p})\),
proving (19).

For \(p=3/2\), split at \(K=\lfloor\mu^{-2}\rfloor\). Replacing
\(1-\cos(2m\mu t)\) by its quadratic term for \(m\le K\) creates total
error \(O(\mu^4K)=O(\mu^2)\), by (11) with \(k=4\). Its main coefficient
is asymptotic to \(\sqrt\pi/m\), whose sum is \(\sqrt\pi\log K\). The
tail is \(O(\mu^2)\) by (21). This proves the first part of (20). The
cubic sine remainder and the first-derivative bound similarly give an
\(O(|\mu|)\) error for the derivative, proving the second part.
\(\square\)

The asymptotics (18)--(19) are approached slowly: the relative
correction is of order \(|\mu|^{3-2p}\), so numerical checks need very
small amplitudes, and near \(p=3/2\) the power law is practically
indistinguishable from the logarithmic case (20). For example, \(p=4/3\)
gives a variance cusp of order \(|\mu|^{5/3}\). These examples directly
refute finite second differentiability for every \(L^2\) observable.
Their leading constants are independent of friction because the
short-time covariance starts as \(1-t^2/2\) for every \(\gamma\).

The regularity is also sharp at first order. For (17) with \(1<p<3/2\),
the derivative is globally Hölder continuous of exponent \(2p-2\), by
summing \(m^{-p}\min(1,\sqrt m|\mu-\nu|)\); (19) rules out a larger
exponent at zero. At \(p=3/2\), this bound gives the small-increment
modulus \(O(\delta\log(1/\delta))\), where \(\delta=|\mu-\nu|\); (20)
rules out Lipschitz continuity at zero. For a single example with no
positive Hölder exponent, take \(\gamma\ge\sqrt2\), \(n_j=2^{2^j}\), and
\(f=\sum_{j\ge1}j^{-1}\phi_{n_j}\). At \(\mu_j=n_j^{-1/2}\), the
derivative local limit gives \(-V_f'(\mu_j)\ge c j^{-2}\) for large
\(j\), since every spectral derivative has the same sign. This exceeds
every constant multiple of \(\mu_j^\alpha\) along the sequence, for
every \(\alpha>0\), although Theorem 5 still gives \(V_f\in C^1\).

\subsection{6. Large amplitudes: polynomial and uniform
rates}\label{large-amplitudes-polynomial-and-uniform-rates}

For any \(\gamma>0\), the Riemann--Lebesgue lemma in (6), dominated
using (5), gives

\[
\lim_{|\mu|\to\infty}V_f(\mu)=V_{\Pi_Jf}(0). \tag{22}
\]

This is DNP's known projection limit. We now quantify it. Define

\[
D_J(\mu)=\sup_{\substack{f\in L^2_0(\pi)\text{ real}\\\|f\|_2=1}}
\{V_f(\mu)-V_{\Pi_Jf}(0)\}
=\sup_{n,\omega\ne0}S_n(\mu\omega). \tag{23}
\]

The last equality follows by concentrating the spectral mass on a
conjugate frequency pair. We define the empty supremum to be zero when
\(J=0\), so \(D_0\equiv0\). The remainder is nonnegative, including when
\(\gamma<\sqrt2\).

\subsubsection{6.1. Every fixed polynomial has fourth-order
decay}\label{every-fixed-polynomial-has-fourth-order-decay}

\textbf{Proposition 7.} For a polynomial observable,

\[
V_f(\mu)-V_{\Pi_Jf}(0)
=\frac{2\gamma}{\mu^4}
\sum_{n,\omega\ne0}\frac{n a_{n,\omega}}{\omega^4}
+O_{f,J,\gamma}(|\mu|^{-6}). \tag{24}
\]

The leading coefficient is positive exactly when \(f\) is non-invariant.

\textbf{Proof.} For fixed \(n\), put \(g_n(t)=r_\gamma(t)^n\) on
\(t\ge0\). Every derivative decays exponentially up to a polynomial
factor. Also \(g_n'(0)=0\) and \(g_n'''(0)=n\gamma\). Six integrations
by parts in the cosine transform give
\(S_n(\xi)=2n\gamma\xi^{-4}+O_{n,\gamma}(|\xi|^{-6})\). A polynomial has
only finitely many chaos components and frequencies, so summing gives
(24). \(\square\)

\subsubsection{6.2. Exact uniform rate for commensurate
rotations}\label{exact-uniform-rate-for-commensurate-rotations}

Let \(G_J\) be the additive subgroup of \(\mathbb R\) generated by the
positive rotation speeds of \(J\). The union of all chaos frequencies
equals \(G_J\): positive and negative eigenvectors in symmetric tensor
powers realise every integer combination. If \(G_J\) is discrete and
\(J\ne0\), write \(G_J=\omega_*\mathbb Z\), where \(\omega_*>0\). This
is exactly the case of commensurate nonzero rotation speeds.

\textbf{Theorem 8.} For every \(\gamma>0\) and every fixed
finite-dimensional \(J\), \(D_J(\mu)\to0\) as \(|\mu|\to\infty\). If
\(G_J=\omega_*\mathbb Z\) with \(\omega_*>0\), then

\[
\boxed{\displaystyle
D_J(\mu)\sim\frac{\sqrt{2\pi/e}}{|\mu|\omega_*}.} \tag{25}
\]

\textbf{Proof.} Lemma 4 gives \(S_n(\xi)\le C_\gamma n^{-1/2}\)
uniformly in \(\xi\). High degrees are therefore uniformly small. Each
bounded set of degrees contains finitely many nonzero frequencies, and
their transforms vanish at infinity. This proves \(D_J\to0\) without any
arithmetic assumption.

To obtain the exact constant, strengthen the scaling in (13). Set
\(g_n(u)=r_\gamma(|u|/\sqrt n)^n\). The proof of Lemma 4 also gives

\[
g_n\longrightarrow e^{-u^2/2}\quad\text{in }W^{1,1}(\mathbb R).
\]

Indeed, near zero the derivatives converge to \(-u e^{-u^2/2}\) and are
bounded by \(C|u|e^{-cu^2}\). Outside \(|u|\le\delta\sqrt n\), their
\(L^1\) norm is at most
\(2n\rho^{n-2}\int_\delta^\infty|rr'|\,\mathrm dt\), which tends to
zero. Fourier transformation of the derivatives yields

\[
\sup_{x\in\mathbb R}\left|x\left[
\sqrt n S_n(\sqrt n x)-\sqrt{2\pi}e^{-x^2/2}
\right]\right|\longrightarrow0. \tag{26}
\]

For \(\omega\ne0\), set \(x=|\mu\omega|/\sqrt n\). The expression
\(|\mu|\omega_*S_n(\mu\omega)\) equals
\((\omega_*/|\omega|)x\sqrt n S_n(\sqrt n x)\). For sufficiently large
\(n\), (26) bounds it by
\(\max_{x\ge0}\sqrt{2\pi}x e^{-x^2/2}+o(1)=\sqrt{2\pi/e}+o(1)\),
uniformly in the frequencies. Every finite set of low degrees
contributes \(o(1)\) after multiplication by \(|\mu|\), by the
fixed-degree tail used in Proposition 7. This proves the upper bound in
(25).

The frequency \(\omega_*\) occurs in some degree \(n_0\) and in every
degree \(n_0+2k\): append one positive and one negative eigenvector from
any rotation plane to a nonzero symmetric tensor. Choose such a degree
nearest to \((\mu\omega_*)^2\). The local limit (13), evaluated at
\(x\to1\), gives the matching lower bound. \(\square\)

\textbf{Corollary 8A (an elementary bound at every amplitude).} For
\(\gamma\ge\sqrt2\) there is also a simple bound at every nonzero
amplitude; unlike (25), it needs the threshold. Every symmetric
decreasing probability density \(h\) satisfies \(h(x)\le1/(2|x|)\), by
integrating over \([-|x|,|x|]\). Applying this to \(S_n/(2\pi)\) gives

\[
0\le V_f(\mu)-V_{\Pi_Jf}(0)
\le\frac{\pi}{|\mu|}
\sum_{n,\omega\ne0}\frac{a_{n,\omega}}{|\omega|}. \tag{27}
\]

When a gap \(\omega_*\) holds, the right side is at most
\(\pi\|f-\Pi_Jf\|_2^2/(|\mu|\omega_*)\). Formula (25) sharpens its
asymptotic constant.

The observable attaining the lower scale in (25) changes with \(\mu\)
and has degree of order \(\mu^2\). Thus (24) and (25) address different
quantifiers: one fixed polynomial versus the worst unit-norm \(L^2\)
observable at each amplitude.

\subsection{7. Arithmetic rates and small
divisors}\label{arithmetic-rates-and-small-divisors}

There need not be a positive angular gap when several rotation speeds
are incommensurate. This changes the uniform rate without destroying
convergence.

\textbf{Proposition 9 (Diophantine bound).} Suppose there exist \(c>0\)
and \(\tau\ge0\) such that every nonzero frequency in degree \(n\)
satisfies \(|\omega|\ge c n^{-\tau}\). Then, for every \(\gamma>0\),

\[
D_J(\mu)=O(|\mu|^{-1/(2\tau+1)}). \tag{28}
\]

\textbf{Proof.} The \(W^{1,1}\) estimate used in (26), together with the
finitely many remaining degrees, gives \(|\xi|S_n(\xi)\le C_\gamma\)
uniformly in \(n,\xi\). Combining it with Lemma 4 gives

\[
S_n(\mu\omega)\le C\min\{n^{-1/2},n^\tau/|\mu|\}.
\]

Split the degrees at \(n\asymp|\mu|^{2/(2\tau+1)}\) to obtain (28).
\(\square\)

\textbf{A sharp irrational example.} Let \(d=4\) and let the two
rotation speeds be \(1\) and \(\sqrt2\). Every degree-\(n\) frequency is
\(a+b\sqrt2\) with integers \(|a|+|b|\le n\). If it is nonzero, then

\[
|a+b\sqrt2|\ge\frac1{|a-b\sqrt2|}\ge\frac1{\sqrt2 n},
\]

because \(a^2-2b^2\) is a nonzero integer. Proposition 9 gives
\(D_J(\mu)=O(|\mu|^{-1/3})\).

For the lower bound, write \(p_k+q_k\sqrt2=(1+\sqrt2)^k\), and put
\(n_k=p_k+q_k\), \(\omega_k=|p_k-q_k\sqrt2|=(1+\sqrt2)^{-k}\). These
frequencies occur in degree \(n_k\). With \(\mu_k=\sqrt{n_k}/\omega_k\),
we have \(\mu_k\asymp n_k^{3/2}\) and bounded successive ratios
\(\mu_{k+1}/\mu_k\). For \(\mu_k\le|\mu|<\mu_{k+1}\), the argument
\(|\mu|\omega_k/\sqrt{n_k}\) remains in a fixed compact subset of
\((0,\infty)\). The uniform local limit (13) then gives
\(S_{n_k}(\mu\omega_k)\ge c n_k^{-1/2}\ge c'|\mu|^{-1/3}\) for large
\(k\). Hence

\[
D_J(\mu)=\Theta(|\mu|^{-1/3}) \tag{29}
\]

for this explicit four-dimensional example.

\textbf{No algebraic rate in general.} Let the two speeds instead be
\(1\) and a positive Liouville number, for example
\(\alpha=\sum_{j\ge1}10^{-j!}\). It admits positive integer approximants
\(p_k/q_k\) with increasing denominators and

\[
0<\omega_k=|q_k\alpha-p_k|<2q_k^{-k},\qquad n_k=p_k+q_k\asymp q_k.
\]

A unit real chaos element \(\psi_k\) of degree \(n_k\) has frequency
pair \(\pm\omega_k\). At \(\mu_k=\sqrt{n_k}/\omega_k\), (13) gives

\[
D_J(\mu_k)\ge S_{n_k}(\sqrt{n_k})
\sim\sqrt{2\pi/e}\,n_k^{-1/2}.
\]

For every \(\beta>0\), the quantity \(\mu_k^\beta D_J(\mu_k)\) diverges.
Therefore \(D_J\) has no bound \(O(|\mu|^{-\beta})\) for any
\(\beta>0\), despite tending to zero. The same obstruction holds for one
fixed observable: after retaining distinct degrees, take
\(f=c\sum_{k\ge1}k^{-1}\psi_k\) with unit \(L^2\) norm. Each spectral
term is nonnegative, so \(V_f(\mu_k)\ge c^2k^{-2}S_{n_k}(\sqrt{n_k})\);
the same comparison proves failure of every algebraic rate. This example
has \(\Pi_Jf=0\).

\subsection{8. A sharp matrix criterion and replica
application}\label{a-sharp-matrix-criterion-and-replica-application}

The scalar covariance in (7) is not essential. We now allow the
stiffness and friction matrices to fail to commute. In this section
\(K\) denotes target precision, rather than the coefficient of a
quadratic observable.

Let \(K,\Gamma\in\mathbb R^{d\times d}\) be symmetric positive definite,
and let \(J^\top=-J\) satisfy \([J,K]=[J,\Gamma]=0\). Consider

\[
\begin{aligned}
\mathrm dq_t&=(p_t-\mu Jq_t)\,\mathrm dt,\\
\mathrm dp_t&=(-Kq_t-\Gamma p_t-\mu Jp_t)\,\mathrm dt
                    +\sqrt{2\Gamma}\,\mathrm dW_t.
\end{aligned} \tag{M1}
\]

The stationary law is \(\mathcal N(0,K^{-1})\otimes\mathcal N(0,I)\).
Define \(V_f\) by (2) and \(\Pi_J\) using rotations of position as
before.

\textbf{Theorem 10 (exact matrix boundary).} Set

\[
\boxed{\displaystyle D=\Gamma^3-\Gamma K-K\Gamma.} \tag{M2}
\]

If \(D\succeq0\), then \(V_f\) is nonincreasing in \(|\mu|\) for every
real square-integrable position observable. Equality between distinct
nonnegative amplitudes holds exactly when \(f=\Pi_Jf\). If \(J\) is
invertible, the following are equivalent:

\begin{enumerate}
\def\labelenumi{\arabic{enumi}.}
\tightlist
\item
  \(D\succeq0\).
\item
  This global ordering holds for every real \(L^2\) position observable.
\item
  Every real linear position observable satisfies \(V_f(\mu)\le V_f(0)\)
  for all sufficiently small \(\mu\).
\end{enumerate}

The matrix assumption \([K,\Gamma]=0\) is absent. Strictness uses
\(-S'(\xi)\succ0\) for \(\xi>0\), which holds even when \(B\) is
singular (Section 8.1). The two-replica construction of Section 8.4
always produces an invertible \(J\), so for replicated targets the three
conditions are equivalent. The symmetry assumptions involving \(J\) do
restrict a single-copy system: an invertible commuting skew matrix
requires even-dimensional eigenspaces of \(K\). Section 8.4 obtains this
symmetry for an arbitrary base Gaussian target by using two replicas.
For singular \(J\), (M2) remains sufficient; no necessity statement is
made for an arbitrary unrotated block.

\subsubsection{8.1. Matrix spectra and tensor
convolution}\label{matrix-spectra-and-tensor-convolution}

At \(\mu=0\), the deterministic equation is
\(\ddot q+\Gamma\dot q+Kq=0\). Its drift is stable. Indeed, an eigenpair
of the second-order pencil gives
\(\lambda^2\|v\|^2+\lambda v^*\Gamma v+v^*Kv=0\); the positive real
coefficients force \(\mathop{\mathrm{Re}}\lambda<0\). The stationary
covariance solves the Lyapunov equation, and the commuting rotation
preserves it and commutes with the drift. All covariance functions
therefore decay exponentially, allowing a polynomial prefactor.

The physical-position transfer spectrum is

\[
S_q(\xi)=2H_\xi\Gamma H_\xi^*,\qquad
H_\xi=(K-\xi^2I+i\xi\Gamma)^{-1}.
\]

The inverse has a useful cancellation. Without assuming that \(K\) and
\(\Gamma\) commute,

\[
\begin{aligned}
S_q(\xi)&=2Q(\xi)^{-1},\\
Q(\xi)&=K\Gamma^{-1}K+\xi^2B+\xi^4\Gamma^{-1},\\
B&=\Gamma-K\Gamma^{-1}-\Gamma^{-1}K,\qquad D=\Gamma B\Gamma.
\end{aligned} \tag{M3}
\]

To check (M3), multiply
\((K-\xi^2I-i\xi\Gamma)\Gamma^{-1}(K-\xi^2I+i\xi\Gamma)\): the imaginary
cross terms cancel. Whitening position by \(x=K^{1/2}q\) gives spectrum
\(S(\xi)=K^{1/2}S_q(\xi)K^{1/2}\) and covariance \(C(t)\). Both are real
symmetric and even, and commute with \(J\). Also \(C(0)=I\),
\(\|C(t)\|\le1\), and \(\int\|C(t)\|\,\mathrm dt<\infty\).

The spectrum \(S\) is decreasing in Loewner order on \([0,\infty)\) if
and only if \(B\succeq0\). Sufficiency follows from
\(Q'(\xi)=2\xi(B+2\xi^2\Gamma^{-1})\) and order reversal under
inversion. Conversely, increasing \(Q\) implies
\((Q(\xi)-Q(0))/\xi^2\succeq0\), whose limit is \(B\). When
\(B\succeq0\), the quartic term makes \(-S'(\xi)\) positive definite for
every \(\xi>0\), including boundary cases where \(B\) is singular.

For spectra acting on separate finite-dimensional spaces, define their
tensor convolution by

\[
(A\star_\otimes B)(x)=\int_{\mathbb R}A(x-y)\otimes B(y)\,\mathrm dy. \tag{M4}
\]

This operation uses tensor products of operators.

\textbf{Lemma 10A (tensor convolution keeps Loewner unimodality).} Let
\(A\) and \(B\) be continuous, even, integrable functions from
\(\mathbb R\) to positive semidefinite matrices on finite-dimensional
spaces \(E\) and \(F\), each nonincreasing in Loewner order on
\([0,\infty)\). Suppose \(B\) is continuously differentiable on
\((0,\infty)\), \(B(r)\to0\) as \(r\to\infty\), and
\(\int_0^\infty\|B'(r)\|\,\mathrm dr<\infty\). Then \(A\star_\otimes B\)
is continuous, even, integrable, positive semidefinite and nonincreasing
in Loewner order on \([0,\infty)\). If, in addition, \(A(a)-A(b)\succ0\)
for \(0\le a<b\) and \(-B'(r)\succ0\) for \(r>0\), then
\((A\star_\otimes B)(y)-(A\star_\otimes B)(x)\succ0\) for \(0\le y<x\).

\textbf{Proof.} Since \(B\) is even and vanishes at infinity,
\(B(y)=\int_{|y|}^\infty[-B'(r)]\,\mathrm dr\) with \(-B'(r)\succeq0\).
Fubini's theorem, justified by integrability of \(\|A\|\) and
\(\|B'\|\), gives the matrix layer-cake formula

\[
(A\star_\otimes B)(x)=
\int_0^\infty H_r(x)\otimes[-B'(r)]\,\mathrm dr,
\qquad H_r(x)=\int_{x-r}^{x+r}A(u)\,\mathrm du .
\]

For \(0\le y<x\),
\(H_r(x)-H_r(y)=\int_{y-r}^{x-r}[A(u+2r)-A(u)]\,\mathrm du\). On this
interval \(u\ge-r\), so \(|u+2r|\ge|u|\), and Loewner monotonicity of
the even function \(A\) gives \(A(u+2r)-A(u)\preceq0\). The tensor
product of a negative semidefinite and a positive semidefinite matrix is
negative semidefinite, because its eigenvalues are products of their
eigenvalues. Integrating over \(r\) proves the weak statement. Under the
strict hypotheses, \(|u+2r|>|u|\) except at the single endpoint
\(u=-r\), so \(H_r(x)-H_r(y)\prec0\) for every \(r>0\); tensoring with
\(-B'(r)\succ0\) and integrating over \(r>0\) gives a negative definite
difference. Positivity, evenness, continuity and integrability are
immediate from the definition. \(\square\)

We apply Lemma 10A to the whitened spectrum \(S\). When \(B\succeq0\),
\(S\) is smooth and even, \(S(\xi)\succ0\), \(-S'(\xi)\succ0\) for
\(\xi>0\), and \(S(\xi)=O(\xi^{-4})\); also
\(\int_0^\infty\|S'(r)\|\,\mathrm dr\le\int_0^\infty\operatorname{tr}[-S'(r)]\,\mathrm dr=\operatorname{tr}S(0)\).
Induction on \(n\), with \(A=S^{\star_\otimes(n-1)}\) and \(B=S\), shows
that every tensor convolution power \(S^{\star_\otimes n}\) is positive
definite and strictly decreasing in Loewner order on \([0,\infty)\). No
common eigenbasis for spectra at different frequencies is used.

Under the Wick isometry from symmetric tensors to Hermite chaos, the
degree-\(n\) covariance operator is \(C(t)^{\otimes n}\) restricted to
the symmetric tensor space. The orthogonal projection
\(\Pi_{\mathrm{sym}}\) onto that space commutes with
\(C(t)^{\otimes n}\), so the degree-\(n\) spectral operator is

\[
\mathcal S_n(\xi)=(2\pi)^{1-n}\,\Pi_{\mathrm{sym}}\,S^{\star_\otimes n}(\xi)\,\Pi_{\mathrm{sym}}\big|_{\mathrm{sym}} .
\]

Compression by a projection preserves Loewner inequalities, and
preserves strict ones on the range of the projection, because
\(\langle v,\Pi X\Pi v\rangle=\langle v,Xv\rangle\) for \(v\) in that
range. Hence \(\mathcal S_n\) is positive definite and strictly
decreasing in Loewner order on \([0,\infty)\) when \(B\succeq0\).

Decompose each chaos into angular frequency components \(f_{n,\omega}\).
These projections commute with the covariance and spectral operators.
Gaussian conditioning gives

\[
V_f(\mu)=\sum_{n\ge1,\omega}
\langle f_{n,\omega},\mathcal S_n(\mu\omega)f_{n,\omega}\rangle. \tag{M5}
\]

Both conjugate frequency signs are included. The bound
\(|\mathop{\mathrm{Cov}}(f(x_t),f(x_0))|\le\|f\|_2^2\|C(t)\|\) justifies
covariance integration, chaos truncation and the variance limit for
every \(L^2\) observable. Formula (M5) proves sufficiency and strictness
in Theorem 10 exactly as (6) did in the scalar setting.

\subsubsection{8.2. Necessity and a noncommuting
witness}\label{necessity-and-a-noncommuting-witness}

If \(J\) is invertible, decompose position space into speed subspaces
\(W_\lambda=\ker(J^2+\lambda^2I)\), \(\lambda>0\). The matrices in (M3)
preserve them. For \(f(q)=a^\top q\) with \(a\in W_\lambda\), symmetry
kills the sine term in the covariance, giving
\(V_f(\mu)=a^\top S_q(\lambda\mu)a\).

If \(B\) is not positive semidefinite, one such subspace contains \(v\)
with \(v^\top Bv<0\). Set \(Q_0=K\Gamma^{-1}K\) and \(a=Q_0v\).
Expansion of the inverse in (M3) gives

\[
V_f''(0)=-4\lambda^2v^\top Bv>0.
\]

Thus even the local linear-observable condition forces \(B\succeq0\),
proving necessity.

For an explicit noncommuting failure, take base matrices

\[
K_0=\begin{pmatrix}2&1\\1&3\end{pmatrix},\qquad
\Gamma_0=\begin{pmatrix}2&0\\0&3\end{pmatrix}.
\]

The diagonal checks \(\gamma_i^2\ge2(K_0)_{ii}\) hold, but
\(D_0=\begin{psmallmatrix}0&-5\\-5&9\end{psmallmatrix}\) is indefinite.
Duplicate the base matrices and rotate the replicas as in Section 8.4.
For \(f(q)=11q^{(1)}_1+13q^{(1)}_2\),

\[
V_f(\mu)=
\frac{1498\mu^4+4312\mu^2+2950}
{\mu^8+3\mu^6+11\mu^4+34\mu^2+25}
=118+12\mu^2+O(\mu^4).
\]

This is an exact witness of increasing variance. For a noncommuting
admissible example, use \(K=\mathop{\mathrm{diag}}(1,4)\otimes I_2\),
\(\Gamma=\begin{psmallmatrix}4&1\\1&4\end{psmallmatrix}\otimes I_2\),
and \(J=I_2\otimes\begin{psmallmatrix}0&-1\\1&0\end{psmallmatrix}\). Its
criterion is
\(D=\begin{psmallmatrix}68&44\\44&44\end{psmallmatrix}\otimes I_2\succ0\);
the base determinant is \(1056\).

Independent isotropic blocks give the useful special case
\(K=\bigoplus_jk_jI\), \(\Gamma=\bigoplus_j\gamma_jI\). If each block
contains a nonzero rotating plane, the exact universal condition is
\(\gamma_j^2\ge2k_j\) for every block. This follows by applying (M2) and
testing a linear observable within each rotating block.

\subsubsection{8.3. Exact friction design
consequences}\label{exact-friction-design-consequences}

For a fixed positive definite shape \(G\) commuting with \(J\), put
\(\Gamma=\gamma G\). The exact threshold is

\[
\gamma^2\ge\lambda_{\max}
\left[G^{-3/2}(GK+KG)G^{-3/2}\right]. \tag{M6}
\]

It follows by congruence from \(D=\gamma(\gamma^2G^3-GK-KG)\). The
largest eigenvalue is positive because the matrix inside the brackets
has trace \(2\operatorname{tr}(KG^{-2})>0\). For scalar friction, (M6)
becomes \(\gamma^2\ge2\lambda_{\max}(K)\).

\textbf{Corollary 11 (least admissible friction).} Among positive
definite matrices satisfying (M2),

\[
\Gamma_*=(2K)^{1/2}
\]

is the unique least element in Loewner order. In particular, it uniquely
minimises \(\operatorname{tr}\Gamma\), with minimum
\(\sqrt2\operatorname{tr}\sqrt K\).

\textbf{Proof.} Put \(H=\Gamma^2/2-K\). The inequality (M2) is the
positive Lyapunov equation \(\Gamma H+H\Gamma=D\succeq0\), so

\[
H=\int_0^\infty e^{-\Gamma t}D e^{-\Gamma t}\,\mathrm dt\succeq0. \tag{M7}
\]

Therefore \(\Gamma^2\succeq2K\). Monotonicity of the positive square
root gives \(\Gamma\succeq\sqrt{2K}\). This standard operator fact also
follows directly from
\(\sqrt A=\pi^{-1}\int_0^\infty t^{-1/2}A(A+tI)^{-1}\,\mathrm dt\) and
order reversal under inversion. Substituting \(\Gamma_*\) gives \(D=0\).
Equality of traces under a semidefinite ordering forces equality of
matrices. Also \(\Gamma_*\) commutes with \(J\) whenever \(K\) does.
\(\square\)

This minimises friction subject to universal rotation ordering: it is
the smallest friction at which rotation is safe for every observable, a
safety boundary rather than a recommendation (compare the remark at the
end of Section 4.1). It does not minimise variance for a particular
observable or numerical computing cost; the observable-specific friction
optimisation of Chak et al.~(2023) has a different objective.

The weaker condition \(\Gamma^2\succeq2K\) is not sufficient. Indeed,
\(\Gamma=\mathop{\mathrm{diag}}(1,2)\) and
\(K=\begin{psmallmatrix}1/4&-1/4\\-1/4&7/4\end{psmallmatrix}\) give
\(\Gamma^2-2K=\tfrac12\begin{psmallmatrix}1&1\\1&1\end{psmallmatrix}\succeq0\),
but \(D=\tfrac14\begin{psmallmatrix}2&3\\3&4\end{psmallmatrix}\) has
determinant \(-1/16\). Duplication again supplies the required rotation.
Thus the admissible matrices do not form the entire upper Loewner cone
above \(\Gamma_*\).

\subsubsection{8.4. Two replicas for an arbitrary Gaussian
target}\label{two-replicas-for-an-arbitrary-gaussian-target}

Let \(K_0,\Gamma_0>0\) be arbitrary symmetric base matrices, with no
common rotation symmetry or mutual commutation assumption. On two copies
put

\[
K=I_2\otimes K_0,\quad \Gamma=I_2\otimes\Gamma_0,\quad
J=\begin{pmatrix}0&-1\\1&0\end{pmatrix}\otimes I. \tag{M8}
\]

Then \(J\) is invertible and commutes with both matrices. The exact
criterion reduces to \(\Gamma_0^3-\Gamma_0K_0-K_0\Gamma_0\succeq0\). At
each stationary time the replicas are independent with the desired
marginal law; their trajectories are coupled by rotation.

Replica coupling and the use of both replicas to estimate expectations
are established ideas (Ohzeki \& Ichiki, 2015; Ichiki \& Ohzeki, 2021).
Kobayashi, Nakagawa and Sasa (2023, supplementary equations (S.71),
(S.75)--(S.76)) explicitly use kinetic duplication and average both
copies, with a different coupling through Hamiltonian terms. Related
generalised auxiliary-variable and finite-force formulations also appear
in Ohzeki (2026). The claims here are the sharp kinetic matrix criterion
and the explicit universal variance consequences for (M8).

Let a centred base observable have Hermite decomposition
\(f=\sum_{n\ge1}f_n\), after whitening. Write \(V_n\) for the
unperturbed base variance of \(f_n\), and \(V_f^0=\sum_nV_n\). Define
\(b_{2m}=\binom{2m}{m}/4^m\) and \(b_{2m+1}=0\).

\textbf{Corollary 12 (replica averages).} Under the matrix criterion,
the variance of the single-copy observable \(f(q^{(1)})\) strictly
decreases for every nonconstant \(f\). For the average of the same two
copies,

\[
F(q^{(1)},q^{(2)})=\tfrac12\bigl(f(q^{(1)})+f(q^{(2)})\bigr),
\]

the variance strictly decreases unless \(f\) belongs to the second
Hermite chaos. A centred quadratic form, with no linear component, gives
constant variance. The exact limiting formulas are

\[
V_F(0)=\tfrac12V_f^0,\qquad
\lim_{|\mu|\to\infty}V_F(\mu)=\sum_{m\ge1}b_{2m}V_{2m}. \tag{M9}
\]

The single-copy observable has the same limiting variance, with baseline
\(V_f^0\). The limits hold for all \(\Gamma_0>0\); the matrix criterion
is needed for monotonicity.

\textbf{Proof.} Write \(c_n(t)\) for the base covariance of \(f_n\). The
first-copy covariance on degree \(n\) is \(c_n(t)\cos^n(\mu t)\). The
two cross-copy covariance matrices have opposite sine signs.
Consequently the covariance of the two-copy average is

\[
\begin{cases}
\tfrac12c_n(t)\cos^n(\mu t),&n\text{ odd},\\
\tfrac12c_n(t)\bigl[\cos^n(\mu t)+\sin^n(\mu t)\bigr],&n\text{ even}.
\end{cases} \tag{M10}
\]

The constant Fourier coefficients of \(\cos^{2m}\) and \(\sin^{2m}\) are
both \(b_{2m}\); odd powers of cosine have zero constant coefficient.
The Riemann--Lebesgue lemma removes all other frequencies on covariance
integration. The integrable chaos bound from (M5) justifies passage to
the limit for general \(L^2\) observables.

A single-copy observable invariant under a quarter-turn would satisfy
\(f(q^{(1)})=f(q^{(2)})\) for independent replicas, so it is constant.
For the average, the degree-\(n\) squared norm is \(\|f_n\|_2^2/2\),
whereas its invariant projection has squared norm \(b_n\|f_n\|_2^2\).
These coincide exactly for \(n=2\); for all other \(n\ge1\),
\(b_n<1/2\). Orthogonality of chaoses and Theorem 10 give the strictness
statement. \(\square\)

Relative to two independent copies, the limiting residual factors for
pure degrees \(2,4,6\) are respectively \(1,3/4,5/8\). Every odd
observable has zero limiting variance. This is an antithetic
cancellation: the rotation carries each replica through its own
negative, so odd components average out. If \(f\) has no second-chaos
component, (M9) is at most \(3/4\) of its two-copy baseline. These
comparisons use continuous-time averages at the same replica count. They
do not establish a gain per unit of numerical work.

Appendix B proves that the exact curvature domain and the sharp
angular-gap constant also hold for the matrix model, even below its
ordering threshold. Thus the extension retains the distinction between
individual polynomial rates and uniform rates over \(L^2\) observables.

\subsection{9. Observation rules and the scope of the
theorem}\label{observation-rules-and-the-scope-of-the-theorem}

\subsubsection{9.1. Phase-space
observables}\label{phase-space-observables}

The position restriction is essential. For a unit momentum observable in
a rotation plane of speed \(\lambda\), the spectrum gives

\[
V_p(\mu)=(\lambda\mu)^2S_\gamma(\lambda\mu).
\]

This is zero at \(\mu=0\) and positive otherwise. At zero rotation, the
time integral of momentum is a difference of position endpoints. Thus
Theorem 1 cannot be extended to every function of \((q,p)\).

\subsubsection{9.2. A fixed finite
horizon}\label{a-fixed-finite-horizon}

Take \(\gamma=2\), the standard planar rotation, \(f(q)=q_1\), and
horizon \(T=1\). Then \(r(t)=(1+t)e^{-t}\) and

\[
W(\mu)=\mathop{\mathrm{Var}}\left(\int_0^1q_1(t)\,\mathrm dt\right)
=2\int_0^1(1-t^2)e^{-t}\cos(\mu t)\,\mathrm dt.
\]

With \(z=-1+i\mu\), direct integration gives

\[
W(\mu)=2\mathop{\mathrm{Re}}\left[-\frac1z+\frac{2e^z}{z^2}
+\frac{2(1-e^z)}{z^3}\right]. \tag{30}
\]

It follows that \(W(2\pi)\simeq0.01932908316\) and
\(W(3\pi)\simeq0.04028633541\). The reversal can also be certified
without floating-point reliance. At \(u=k\pi\),

\[
W(k\pi)=
\frac{2u^4+16u^2-2+(4/e)(-1)^k(2-3u^2-u^4)}{(1+u^2)^3}.
\]

For \(k=3\) the factor \((-1)^k(2-3u^2-u^4)=u^4+3u^2-2\) is positive,
and for \(k=2\) the factor \(2-3u^2-u^4\) is negative. Using \(4/e>4/3\)
(from \(e<3\)) in each case gives

\[
W(3\pi)>\frac{\tfrac{10}3u^4+20u^2-\tfrac{14}3}{(1+u^2)^3}\Big|_{u^2=9\pi^2},
\qquad
W(2\pi)<\frac{\tfrac23u^4+12u^2+\tfrac23}{(1+u^2)^3}\Big|_{u^2=4\pi^2}.
\]

Each numerator increases and each denominator increases with \(u^2\).
With \(9<\pi^2<10\), bound the numerator at \(u^2=81\) and the
denominator at \(u^2=90\) for \(W(3\pi)\), and the reverse (\(u^2=40\)
and \(u^2=36\)) for \(W(2\pi)\); the right-hand fraction rounds the
resulting bound \(4642/151959\) outward to a simpler one:

\[
W(3\pi)>\frac{70456}{3\cdot91^3}=\frac{70456}{2260713}
>\frac{1548}{50653}=\frac{1548}{37^3}>W(2\pi).
\]

The difference between the middle fractions is
\(69224044/114511895589>0\). In contrast, the long-run variance is
\(4/(1+\mu^2)^2\), strictly decreasing. A long-run ordering therefore
does not establish ordering at each observation horizon.

\subsubsection{9.3. Regular observation and exact
aliasing}\label{regular-observation-and-exact-aliasing}

For the same example, sample at interval \(\Delta>0\). The asymptotic
variance per recorded sample is

\[
\tau_\Delta^2(\mu)=1+2\sum_{k\ge1}(1+k\Delta)e^{-k\Delta}
\cos(k\Delta\mu).
\]

Putting \(z=e^{-\Delta+i\Delta\mu}\) yields

\[
\tau_\Delta^2(\mu)=1+2\mathop{\mathrm{Re}}\left[
\frac z{1-z}+\frac{\Delta z}{(1-z)^2}\right]. \tag{31}
\]

This is a nonconstant periodic function of \(\mu\), of period
\(2\pi/\Delta\), and cannot be globally nonincreasing. In fact, the
entire sampled Gaussian law repeats with this period. This is
observation aliasing even under exact simulation, rather than an error
caused by a numerical integrator. Multiplication by \(\Delta\) converts
(31) to variance per physical observation time.

\subsubsection{9.4. Poisson observation preserves the
ordering}\label{poisson-observation-preserves-the-ordering}

Observe at arrival times of an independent Poisson process of rate
\(\lambda>0\), and average a fixed number \(N\) of observations. Let
\(\tau_{\mathrm P,\lambda}^2(\mu)\) denote the limit of \(N\) times the
variance of this sample mean, averaging over both the stationary
dynamics and the independent observation clock. The gap between
observations \(k\) positions apart has the Gamma\((k,\lambda)\) density.
Since

\[
\sum_{k\ge1}\frac{\lambda^kt^{k-1}e^{-\lambda t}}{(k-1)!}=\lambda,
\]

integrable covariance and the covariance-sum identity give

\[
\boxed{\displaystyle
\tau_{\mathrm P,\lambda}^2(\mu)=\|f\|_2^2+\lambda V_f(\mu).} \tag{32}
\]

Thus independent Poisson observation preserves every ordering in Theorem
1 at every observation rate. This is an instance of the classical fact
that Poisson sampling is free of aliasing (Shapiro \& Silverman, 1960).
Per mean elapsed time, the corresponding variance is
\(V_f(\mu)+\|f\|_2^2/\lambda\). This extra term is an observation noise
floor, so (32) does not establish superiority over regularly spaced
observations at equal cost. The identity is a direct consequence of
standard Poisson sampling, presented as an application of the theorem
rather than a separate priority claim.

\subsection{10. Reproducibility and
interpretation}\label{reproducibility-and-interpretation}

The archive accompanying this manuscript contains three deterministic
verification scripts: \texttt{code/verify.py},
\texttt{code/verify\_matrix\_extension.py} and
\texttt{code/verify\_below\_threshold.py}. Each writes a JSON report to
a path chosen by the user and fails with an explicit error, also under
optimised Python, if a check fails. They check the covariance reduction
against matrix exponentials, the linear and quadratic variance formulas
against independent linear-system calculations, Fourier transforms
against quadrature, and the finite-horizon and observation-rule
examples. They also record exact polynomial fingerprints and the
high-degree scaling in (13). The matrix script checks the inverse
spectrum against the complex transfer function, tensor convolution
against a quadratic Lyapunov calculation, and both a noncommuting
admissible example and an exact criterion failure. The third script
computes the crossing points of Table 1 from an exact exponential-sum
formula and cross-checks it by quadrature. The checks support the
derivation and catch normalisation errors; the proofs above establish
the infinite-dimensional and all-parameter assertions.

\begin{figure}
\centering
\includegraphics[width=1\linewidth,height=\textheight,keepaspectratio,alt={Exact model calculations. The friction threshold governs the continuous-time long-run variance; higher chaoses inherit the ordering. Fixed horizons and regular observations can reverse it. Independent Poisson observation retains the ordering with an additional observation variance term. These curves illustrate proved formulas and are not Monte Carlo estimates.}]{figures/variance_landscape.pdf}
\caption{Exact model calculations. The friction threshold governs the
continuous-time long-run variance; higher chaoses inherit the ordering.
Fixed horizons and regular observations can reverse it. Independent
Poisson observation retains the ordering with an additional observation
variance term. These curves illustrate proved formulas and are not Monte
Carlo estimates.}
\end{figure}

The scalar theorem settles a stated question about Gaussian kinetic
sampling, including every square-integrable position observable and
every perturbation amplitude. The matrix theorem gives an exact
extension and a friction design rule, and replication makes its symmetry
available for arbitrary base Gaussian targets. The regularity result
identifies a failure in the literal unrestricted derivative formulation.
The arithmetic results explain why a fast decay for each fixed
polynomial does not control the worst observable uniformly.

These conclusions depend on a Gaussian joint law and a
symmetry-preserving matched rotation. The matrix results allow
covariance operators with no common eigenbasis across frequencies. A
non-Gaussian extension would require control of couplings between chaos
degrees or another replacement for (M5). A numerical-efficiency result
would require a specified discretisation, cost model and error
tolerance. The literature search for this work was last updated on 4
November 2026. It covered the originating paper and thesis, the
overdamped variance-reduction literature, work citing DNP, related
nonreversible-sampling work, and formula-level aliases across Gaussian
spectral analysis and convolution theory; the two closest recent
sources, Maini (2026) and Franke et al.~(2021), were read in full. No
earlier statement of the main results was found.

The manuscript has undergone AI-assisted proof checks and computational
verification, but has not undergone external human peer review or formal
proof-assistant verification. Two natural next steps are left open: an
exact discrete-time criterion for a specified splitting integrator with
rotation, where Section 9.3 shows that aliasing is the obstacle, and a
perturbative statement for near-Gaussian targets.

\subsection{References}\label{references}

\subsubsection{Sampling and the originating
question}\label{sampling-and-the-originating-question}

Chak, M., Kantas, N., Lelièvre, T., \& Pavliotis, G. A. (2023). Optimal
friction matrix for underdamped Langevin sampling. \emph{ESAIM:
Mathematical Modelling and Numerical Analysis, 57}(6), 3335--3371.
\url{https://doi.org/10.1051/m2an/2023083}

Duncan, A. B., Lelièvre, T., \& Pavliotis, G. A. (2016). Variance
reduction using nonreversible Langevin samplers. \emph{Journal of
Statistical Physics, 163}, 457--491.
\url{https://doi.org/10.1007/s10955-016-1491-2}

Duncan, A. B., Nüsken, N., \& Pavliotis, G. A. (2017). Using perturbed
underdamped Langevin dynamics to efficiently sample from probability
distributions. \emph{Journal of Statistical Physics, 169}, 1098--1131.
\url{https://doi.org/10.1007/s10955-017-1906-8}

Guillin, A., \& Monmarché, P. (2016). Optimal linear drift for the speed
of convergence of an hypoelliptic diffusion. \emph{Electronic
Communications in Probability, 21}, 1--14.
\url{https://doi.org/10.1214/16-ECP25}

Hwang, C.-R., Hwang-Ma, S.-Y., \& Sheu, S.-J. (1993). Accelerating
Gaussian diffusions. \emph{The Annals of Applied Probability, 3}(3),
897--913. \url{https://doi.org/10.1214/aoap/1177005371}

Hwang, C.-R., Normand, R., \& Wu, S.-J. (2015). Variance reduction for
diffusions. \emph{Stochastic Processes and their Applications, 125}(9),
3522--3540. \url{https://doi.org/10.1016/j.spa.2015.03.006}

Kaiser, M., Jack, R. L., \& Zimmer, J. (2017). Acceleration of
convergence to equilibrium in Markov chains by breaking detailed
balance. \emph{Journal of Statistical Physics, 168}(2), 259--287.
\url{https://doi.org/10.1007/s10955-017-1805-z}

Lelièvre, T., Nier, F., \& Pavliotis, G. A. (2013). Optimal
non-reversible linear drift for the convergence to equilibrium of a
diffusion. \emph{Journal of Statistical Physics, 152}(2), 237--274.
\url{https://doi.org/10.1007/s10955-013-0769-x}

Rey-Bellet, L., \& Spiliopoulos, K. (2015a). Irreversible Langevin
samplers and variance reduction: a large deviations approach.
\emph{Nonlinearity, 28}(7), 2081--2103.
\url{https://doi.org/10.1088/0951-7715/28/7/2081}

Rey-Bellet, L., \& Spiliopoulos, K. (2015b). Variance reduction for
irreversible Langevin samplers and diffusion on graphs. \emph{Electronic
Communications in Probability, 20}, 1--16.
\url{https://doi.org/10.1214/ECP.v20-3855}

Rey-Bellet, L., \& Spiliopoulos, K. (2016). Improving the convergence of
reversible samplers. \emph{Journal of Statistical Physics, 164}(3),
472--494. \url{https://doi.org/10.1007/s10955-016-1565-1}

Franke, B., Hwang, C.-R., Ouled Said, A., \& Pai, H.-M. (2021). A note
on the asymptotic variance of drift accelerated diffusions.
\emph{Statistics \& Probability Letters, 175}, 109128.
\url{https://doi.org/10.1016/j.spl.2021.109128}

Nüsken, N. (2018). \emph{Topics in sampling schemes based on Markov
processes} {[}Doctoral thesis, Imperial College London{]}. Spiral.
\url{https://doi.org/10.25560/63868}

\subsubsection{Replica coupling and related
constructions}\label{replica-coupling-and-related-constructions}

Ichiki, A., \& Ohzeki, M. (2021). \emph{Stochastic dynamics without
detailed balance condition connecting simple gradient method and
Hamiltonian Monte Carlo}. arXiv.
\url{https://doi.org/10.48550/arXiv.2107.08633}

Kobayashi, M., Nakagawa, N., \& Sasa, S.-i. (2023). Control of
metastable states by heat flux in the Hamiltonian Potts model.
\emph{Physical Review Letters, 130}, 247102.
\url{https://doi.org/10.1103/PhysRevLett.130.247102}. Supplementary
material inspected in \url{https://arxiv.org/pdf/2212.12289}.

Ohzeki, M., \& Ichiki, A. (2015). Langevin dynamics neglecting detailed
balance condition. \emph{Physical Review E, 92}, 012105.
\url{https://doi.org/10.1103/PhysRevE.92.012105}

Ohzeki, M. (2026). Nonreversible gauge fields in Fokker-Planck dynamics:
Supersymmetric Hamiltonians and learned finite forces. \emph{Physical
Review Research, 8}, 033170. \url{https://doi.org/10.1103/tg4c-7cnr}

\subsubsection{Spectral and convolution
ingredients}\label{spectral-and-convolution-ingredients}

Berg, C., \& Vignat, C. (2008). Linearization coefficients of Bessel
polynomials and properties of Student t-distributions.
\emph{Constructive Approximation, 27}, 15--32.
\url{https://doi.org/10.1007/s00365-006-0643-6}

Breuer, P., \& Major, P. (1983). Central limit theorems for non-linear
functionals of Gaussian fields. \emph{Journal of Multivariate Analysis,
13}(3), 425--441. \url{https://doi.org/10.1016/0047-259X(83)90019-2}

Maini, L. (2026). \emph{Spectral delocalization and central limit
theorems for Fourier functionals of stationary Gaussian fields}
(arXiv:2610.37113). arXiv.
\url{https://doi.org/10.48550/arXiv.2610.37113}

Lefèvre, C., \& Utev, S. (2013). Convolution property and exponential
bounds for symmetric monotone densities. \emph{ESAIM: Probability and
Statistics, 17}, 605--613. \url{https://doi.org/10.1051/ps/2012012}

Muscinelli, S. P., Gerstner, W., \& Schwalger, T. (2019). How single
neuron properties shape chaotic dynamics and signal transmission in
random neural networks. \emph{PLOS Computational Biology, 15}(6),
e1007122. \url{https://doi.org/10.1371/journal.pcbi.1007122}

\subsubsection{Sampling schedules}\label{sampling-schedules}

Shapiro, H. S., \& Silverman, R. A. (1960). Alias-free sampling of
random noise. \emph{Journal of the Society for Industrial and Applied
Mathematics, 8}(2), 225--248. \url{https://doi.org/10.1137/0108013}

\subsubsection{Related arithmetic variance
phenomena}\label{related-arithmetic-variance-phenomena}

Lachièze-Rey, R. (2022). \emph{Diophantine Gaussian excursions and
random walks} (arXiv:2104.07290, version 4). arXiv.
\url{https://doi.org/10.48550/arXiv.2104.07290}

\newpage

\subsection{Appendix A. Explicit fingerprints and
checks}\label{appendix-a.-explicit-fingerprints-and-checks}

At \(\gamma=2\), the oscillator covariance is \(r(t)=(1+t)e^{-t}\) for
\(t\ge0\). Hence

\[
S_n(\xi)=2\mathop{\mathrm{Re}}\sum_{k=0}^n\binom nk\frac{k!}{(n-i\xi)^{k+1}}. \tag{33}
\]

The first zero-frequency values are \(S_1(0)=4\), \(S_2(0)=5/2\), and
\(S_3(0)=52/27\). The first spectral density is \(4/(1+\xi^2)^2\);
normalised by \(2\pi\), it is a scaled Student distribution with three
degrees of freedom. Berg and Vignat (2008) study the associated
convolution family. These fingerprints locate classical distributional
and spectral ingredients. The separate novelty report also records the
exact connection between \(nS_n(0)/2\) and the integral defining OEIS
sequence A090878.

For \(f=q_1\) and a unit planar rotation, (10) gives the following exact
checks:

\[
\begin{array}{c|ccc}
&\mu=0&\mu=1/2&\mu=1\\\hline
\gamma=1&2&32/13&2\\
\gamma=\sqrt2&2\sqrt2&32\sqrt2/17&\sqrt2\\
\gamma=2&4&64/25&1
\end{array}
\]

The below-threshold row initially increases. The two other rows
decrease, including the threshold itself. At zero perturbation, (10)
gives \(V_{q_1}(0)=2\gamma\); integrating (7) directly gives
\(\int_0^\infty r_\gamma(t)\,\mathrm dt=\gamma\), confirming the factor
of two in (2).

For a real quadratic observable \(f(q)=q^\top Kq-\operatorname{tr}K\),
\(K=K^\top\), let \(A=\mathop{\mathrm{diag}}(K,0)\) and solve

\[
B_\mu^\top C+CB_\mu=A.
\]

Then \(x^\top Cx-\operatorname{tr}C\) solves the Poisson equation for
\(f\), and Wick's identity gives \(V_f(\mu)=4\operatorname{tr}(AC)\).
This is an independent finite-dimensional check on the chaos expansion.
It is not used to infer the result for general \(L^2\) observables.

\newpage

\subsection{Appendix B. Regularity and rates for matrix kinetic
systems}\label{appendix-b.-regularity-and-rates-for-matrix-kinetic-systems}

This appendix extends the regularity and rate arguments to the matrix
model of Section 8. Let \(K,\Gamma>0\) be real symmetric matrices, let
\(J\) be real skew with \([J,K]=[J,\Gamma]=0\), and use stationary
position law \(\mathcal N(0,K^{-1})\). The results here require no
friction-threshold inequality. A kernel of \(J\) is allowed. Coordinates
are whitened by \(x=K^{1/2}q\), so the Hermite decomposition and its
norm refer to standard Gaussian measure. Subtract the mean of \(f\); all
chaos sums below start at \(n=1\).

Write \(C(t)=\mathop{\mathrm{Cov}}(x_t^0,x_0^0)\). It is real symmetric
and even. Its short-time expansion on \(t\ge0\) is

\[
C(t)=I-\tfrac12Kt^2+\tfrac16Lt^3+O(t^4),
\qquad L=K^{1/2}\Gamma K^{1/2}. \tag{B1}
\]

All its derivatives on the positive half-line decay exponentially. Also
\(\|C(t)\|<1\) for every \(t>0\): the conditional covariance
\(I-C(t)C(t)^\top\) contains the strictly positive covariance of the
noise reaching position over \((0,t)\). The position response to
momentum noise is \(sI+O(s^2)\) near zero, and \(\Gamma>0\), which
proves strict positivity in every direction. Consequently, for some
\(c,C,\delta>0\),

\[
e^{-Ct^2}I\preceq C(t)\preceq e^{-ct^2}I
\quad(0\le t\le\delta),\qquad
\rho=\sup_{|t|\ge\delta}\|C(t)\|<1. \tag{B2}
\]

Let \(\mathcal C_n(t)=C(t)^{\otimes n}\) restricted to the symmetric
tensor space, and let \(\mathcal S_n=\widehat{\mathcal C_n}\). Use the
isometric identification of that space with the \(n\)th Hermite chaos.
These spectral operators are positive semidefinite and commute with the
angular projections. If \(f_{n,\omega}\) denotes a projected chaos
component, then

\[
V_f(\mu)=\sum_{n,\omega}
\langle f_{n,\omega},\mathcal S_n(\mu\omega)f_{n,\omega}\rangle. \tag{B3}
\]

The convention includes both signs of each nonzero frequency. The bound
\(\|\mathcal C_n(t)\|\le\|C(t)\|^n\le\|C(t)\|\) makes the sum absolutely
integrable for every \(f\in L^2\).

\subsubsection{B.1. Exact regularity}\label{b.1.-exact-regularity}

\textbf{Proposition B1.} Every real square-integrable position
observable has \(V_f\in C^1(\mathbb R)\) and \(V_f'(0)=0\). It has a
finite ordinary second derivative at zero if and only if

\[
\sum_{n,\omega}\omega^2n^{-3/2}\|f_{n,\omega}\|_2^2<\infty. \tag{B4}
\]

This is also equivalent to \(V_f\in C^2(\mathbb R)\). If (B4) diverges,
\((V_f(\mu)-V_f(0))/\mu^2\to-\infty\) as \(\mu\to0\).

\textbf{Proof.} The split in (B2) gives

\[
\int_0^\infty t^k\|C(t)\|^n\,\mathrm dt\le C_k n^{-(k+1)/2}.
\]

Since \(|\omega|\le n\|J\|\), the first derivative of every matrix term
in (B3) is bounded by a constant times \(\|f_{n,\omega}\|_2^2\). The
weighted second derivative is bounded by a constant times the summand in
(B4). Uniform convergence proves the stated sufficiency and first
differentiability.

For necessity, restrict the cosine-difference integral to
\([0,\delta]\). Its kernels are positive semidefinite by (B2), and

\[
\int_0^\delta t^2\langle v,\mathcal C_n(t)v\rangle\,\mathrm dt
\ge c_1 n^{-3/2}\|v\|^2.
\]

Fatou's lemma therefore forces the positive small-time contribution to
the divided difference to diverge when (B4) diverges. The absolute value
of the complementary contribution is bounded by
\(C_2\|J\|^2\sup_n(n^2\rho^{n-1})\|f\|_2^2\), using
\(1-\cos u\le u^2/2\). It cannot cancel the divergence. This also
excludes a finite ordinary second derivative. \(\square\)

\subsubsection{B.2. A uniform matrix Fourier
approximation}\label{b.2.-a-uniform-matrix-fourier-approximation}

Define the positive operator on the symmetric tensor space

\[
A_n=\sum_{j=1}^n I^{\otimes(j-1)}\otimes K\otimes I^{\otimes(n-j)},
\qquad T_n=A_n/n.
\]

Its spectrum lies in the fixed interval
\([\lambda_{\min}(K),\lambda_{\max}(K)]\). We claim

\[
\sup_{x\in\mathbb R}(1+|x|)\left\|
\sqrt n\,\mathcal S_n(\sqrt n x)
-\sqrt{2\pi}\,T_n^{-1/2}e^{-x^2T_n^{-1}/2}
\right\|\longrightarrow0. \tag{B5}
\]

Here the comparison operator depends on \(n\); no convergence of \(T_n\)
is asserted.

To prove (B5), compare \(C(|u|/\sqrt n)^{\otimes n}\) with
\(e^{-A_nu^2/(2n)}\). On \(|u|\le\delta\sqrt n\), telescope the tensor
factors and use (B1). The norm of the difference is bounded by
\(C_3n^{-1/2}|u|^3e^{-c_3u^2}\). Differentiating the tensor products
adds a bound \(C_4n^{-1/2}(|u|^2+|u|^4)e^{-c_4u^2}\) for the derivative
difference. These bounds follow from \(C(t)-e^{-Kt^2/2}=O(t^3)\) and
\(C'(t)+(Kt)e^{-Kt^2/2}=O(t^2)\); the other tensor factors supply the
Gaussian bound in (B2).

On the complement, for \(n\ge2\), the function and derivative integrals
for the covariance tensor are bounded respectively by

\[
2\sqrt n\,\rho^{n-1}\int_\delta^\infty\|C(t)\|\,\mathrm dt,
\qquad
2n\rho^{n-2}\int_\delta^\infty\|C(t)\|\,\|C'(t)\|\,\mathrm dt.
\]

Both tend to zero; the comparison Gaussian has uniformly Gaussian tails.
Thus the difference tends to zero in operator-norm \(W^{1,1}\). Fourier
transformation, and integration by parts for the frequency-weighted
transform, prove (B5). Symmetric-tensor restriction can only decrease
the operator norm.

\subsubsection{B.3. Exact angular-gap constant and polynomial
asymptotics}\label{b.3.-exact-angular-gap-constant-and-polynomial-asymptotics}

Define \(D_J\) as in (23), using the present Gaussian law and
projection. It equals

\[
\sup_{n,\omega\ne0}\bigl\|\mathcal S_n(\mu\omega)|_{\mathcal H_{n,\omega}}\bigr\|.
\]

Real observables realise the same supremum by combining conjugate
components.

\textbf{Proposition B2.} For any fixed \(K,\Gamma>0\) and commuting skew
\(J\), \(D_J(\mu)\to0\). If the nonzero angular frequency group is
\(\omega_*\mathbb Z\), then

\[
D_J(\mu)\sim\frac{\sqrt{2\pi/e}}{|\mu|\omega_*}. \tag{B6}
\]

The leading constant is independent of \(K\) and \(\Gamma\). The
Diophantine condition in Proposition 9 still gives
\(D_J(\mu)=O(|\mu|^{-1/(2\tau+1)})\).

\textbf{Proof.} The moment estimate gives
\(\|\mathcal S_n\|_\infty\le Cn^{-1/2}\). This controls high degrees,
and finitely many low-degree spectral matrices vanish at infinity. Thus
\(D_J\to0\).

For the upper bound in (B6), put \(x=|\mu\omega|/\sqrt n\) in (B5). For
every positive scalar \(\kappa\),

\[
\sup_{x\ge0}x\kappa^{-1/2}e^{-x^2/(2\kappa)}=e^{-1/2}.
\]

Functional calculus for \(T_n\) therefore bounds the high-degree part of
\(|\mu|\omega_*D_J(\mu)\) by \(\sqrt{2\pi/e}+o(1)\). For each bounded
degree, \(|\xi|\mathcal S_n(\xi)\to0\), by the Riemann--Lebesgue lemma
applied to its integrable derivative. These degrees contribute \(o(1)\).

For the lower bound, \(K\) and \(J\) have a common complex orthonormal
eigenbasis. Choose a nonzero symmetric tensor that has angular frequency
\(\omega_*\) and \(A_{n_0}\) eigenvalue \(a_0>0\). Append \(k\) copies
of a conjugate eigenvector pair. Its angular frequency remains
\(\omega_*\), its degree is \(n_0+2k\), and its stiffness eigenvalue is
\(a_0+2k\kappa\), for some \(\kappa>0\). Choose \(k\) so that this last
value is nearest to \((\mu\omega_*)^2\). The unweighted part of (B5)
then gives the lower constant \(\sqrt{2\pi/e}\). The tensor is nonzero
because it corresponds to a product of nonzero linear polynomials.

Finally, (B5) gives \(\sup_{n,\xi}|\xi|\|\mathcal S_n(\xi)\|<\infty\).
Combining this with \(Cn^{-1/2}\) gives the same degree split as
Proposition 9. \(\square\)

For completeness, fixed polynomials have the matrix counterpart of (24).
Let
\(L_n=\sum_{j=1}^n I^{\otimes(j-1)}\otimes L\otimes I^{\otimes(n-j)}\).
Six integrations by parts, using \(\mathcal C_n'(0)=0\) and
\(\mathcal C_n'''(0)=L_n\), give

\[
V_f(\mu)-V_{\Pi_Jf}(0)
=\frac2{\mu^4}\sum_{n,\omega\ne0}
\frac{\langle f_{n,\omega},L_nf_{n,\omega}\rangle}{\omega^4}
+O_{f,K,\Gamma,J}(|\mu|^{-6}). \tag{B7}
\]

The leading coefficient is strictly positive exactly for non-invariant
\(f\), since every \(L_n\) is positive definite. The scalar formulas are
recovered by \(K=I\), \(\Gamma=\gamma I\).

\end{document}
